%% 第31回編集委員会（物理第5巻）：組 原W1
%% p5-01.tex  原-1　原子の中をのぞく：電子と原子核の発見（★1〜4・3段にしない：N-1）
%% 中身は計画書 tools/phys5-plan.md の B「原-1」。公式の持ち主は A（A2・A3・A6・A15・A16・A26・A43）と A-7（ミリカンの終端速度・質量分析）。
%% 前の巻の約束（計画書 A-6）：Q44（電子の発見 1897年：①②）、Q51（電-6 例題8 の 1.0 kV・1.0 mT・半径 0.1066 m から比電荷：例題3）、
%%   Q62（比電荷：②）。力-2 例題12・力-15（ストークスの抵抗、空気の粘度 1.81×10⁻⁵ Pa·s）は④・例題6で同じ数を使った。
%%   電-4 例題（ファラデー定数 96485 C/mol）は⑦・例題7。電-1 の k＝8.99×10⁹、電-2 ⑥の電子ボルトを使った。
%% 数学との重なり（N-10・W-6）：2-7（10 の累乗）、4-6⑤（tan の逆数としての cot）、4-9④（半角）、8-8⑬（立体角：言葉だけ）を指した。
%% 数値はすべて Python で計算した：checks/物理第5巻/原-1/calc_p5_01.py（出力 calc_p5_01_out.txt）。
%%   ⑧の囲みのプログラムは checks/物理第5巻/原-1/sim_alpha.py（紙面と同じ。出力も同じ）、図の軌道の座標は fig_p5_01.py。
%%   常用漢字の確かめ：checks/物理第5巻/原-1/kanji_p5_01_02.py。
%% 記号（この章）：電気素量 e、電子の質量 m（m_e と同じ。⑤からは α 粒子の質量も m）、比電荷 e/m、電場 E（この章だけ。運動エネルギーは K）、
%%   磁場 B、速さ v、円の半径 r、油滴の半径 a（円の半径 r とぶつけない）、油滴の密度 ρ、空気の粘度 η（力-2・力-15 の「粘度」。
%%   第1巻の摩擦係数 μ・この巻の減弱係数 μ とぶつけない）、極板の間かく d_p（いちばん近づく距離 d とぶつけない）、電圧 V、
%%   原子番号 Z、いちばん近づく距離 d、ねらいの距離 b、散乱角 θ、α 線の運動エネルギー K、クーロンの定数 k（電-1。波数は出さない）、
%%   重力加速度 g、アボガドロ定数 N_A、ファラデー定数は文字にしない（電-4・熱-10 と同じ）、数密度 n、はくの厚さ t_f（時刻 t とぶつけない）。
%% 出典（11月11日に確かめた）
%%   ・定数：CODATA 2022（NIST https://physics.nist.gov/cuu/Constants/Table/allascii.txt を curl で読んだ）：e＝1.602176634×10⁻¹⁹ C（定義値）、
%%       mₑ＝9.1093837139×10⁻³¹ kg、電子の比電荷 1.75882000838×10¹¹ C/kg、N_A＝6.02214076×10²³ /mol（定義値）、F＝96485.33212 C/mol、
%%       ボーア半径 5.29177210544×10⁻¹¹ m、陽子の電荷の半径（rms）0.84075 fm、微細構造定数と ħc から ke²＝1.440 MeV·fm、u＝1.66053906892×10⁻²⁷ kg。
%%   ・原子の質量：NIST "Atomic Weights and Isotopic Compositions"（計画書 D5）：金 196.966570、¹³C 13.00335483507、³⁵Cl 34.968852682、³⁷Cl 36.965902602。
%%   ・金の密度 19.32 g/cm³・鉄 7.874 g/cm³：NIST "X-Ray Mass Attenuation Coefficients" 表1（計画書 D7）。金の原子の半径 0.144 nm は、
%%       この密度と面心立方（格子に4個）から計算した（格子定数 0.408 nm、となりの原子の中心の間の半分。calc_p5_01.py）。
%%   ・原子核の半径の目安 R≈1.2 A^{1/3} fm：原子核物理の教科書の標準の目安（K. S. Krane "Introductory Nuclear Physics" 3章ほか）。詳しくは原-11。
%%   ・ノーベル賞：トムソン 1906年物理学賞、ミリカン 1923年物理学賞、ラザフォード 1908年化学賞、チャドウィック 1935年物理学賞（計画書 D13・D18・D40・D23、nobelprize.org）。
%%   ・トムソンの論文：J. J. Thomson, "Cathode Rays", Philosophical Magazine 44, 293 (1897年10月)。ミリカン：R. A. Millikan, "On the Elementary Electrical
%%       Charge and the Avogadro Constant", Phys. Rev. 2, 109 (1913)（APS News "This Month in Physics History" 2006年8月・AIP の展示で照らし合わせた）。
%%       ラザフォード：Philosophical Magazine 21, 669 (1911年5月)（St Andrews の C. Baily "Atomic models" EPJ H (2013) の本文の書誌で確かめた）。
%%       ガイガーとマースデンの実験 1909年、チャドウィックの中性子 1932年は、上のノーベル賞の発表文と Baily (2013) で照らし合わせた。
%%   ・陽子半径の食いちがい：R. Pohl ほか Nature 466, 213 (2010)（ミュー粒子の水素：0.84 fm）。W. Xiong ほか（PRad）"A small proton charge radius from
%%       an electron–proton scattering experiment", Nature 575, 147 (2019)、0.831±0.007±0.012 fm（ADS の要旨と Chemistry World の記事で照らし合わせた）。
%%       CODATA 2022 は 0.84075 fm。本文は「小さいほうの値にまとまりつつあるが、測定と検討は続いている」と書いた。
%%   ・空気の粘度 1.81×10⁻⁵ Pa·s（20 ℃）：第1巻 力-2 例題12 と同じ値（巻末の表）。
%% 〔要確認〕なし。例題の仮の値（油滴の密度 870 kg/m³、はくの厚さ 0.40 μm、トムソンの装置の電場と磁場、電気分解の電流と時間）は本文に「〜とします」と書いた。
\makeatletter
% 原子核（＋の丸）・電子（−の丸）・α 粒子  \txbnuc{(中心)}{半径}  \txbel{(中心)}{半径}
\newcommand\txbnuc[2]{\filldraw[fill=bkRed!75!black, draw=bkInk, line width=0.4pt] #1 circle (#2);}
\newcommand\txbel[2]{\filldraw[fill=bkBlue!70!black, draw=bkInk, line width=0.4pt] #1 circle (#2);\node[font=\sffamily\bfseries\scriptsize, text=white, inner sep=0pt] at #1 {$-$};}
% 紙のおくへ向かう磁場（⊗）  \txbin{(中心)}{半径}
\newcommand\txbin[2]{\begin{scope}[shift={#1}]\draw[bkBlue!70!black, line width=0.5pt] (0,0) circle (#2); \draw[bkBlue!70!black, line width=0.5pt] (45:#2) -- (225:#2) (135:#2) -- (315:#2);\end{scope}}
\tikzset{txb lab/.style={font=\sffamily\footnotesize, text=bkInk},
  txb beam/.style={bkOrange!85!black, line width=1.3pt},
  txb F/.style={bkRed!80!black, line width=1.3pt, -{Stealth[length=2.4mm]}},
  txb Fb/.style={bkBlue!75!black, line width=1.3pt, dash pattern=on 3pt off 1.5pt, -{Stealth[length=2.4mm]}},
  txb dim/.style={bkGray, line width=0.5pt, {Stealth[length=1.6mm]}-{Stealth[length=1.6mm]}},
  txb plate/.style={fill=bkGray!45, draw=bkInk, line width=0.5pt},
  txb path/.style={bkPurple!80!black, line width=0.9pt}}
\makeatother

\phchapter{1}{原子の中をのぞく：\\電子と原子核の発見}{原子の中は、ほとんどからっぽだって本当？}{\bklevelrange{1}{4}}{原子の中をのぞく：電子と原子核の発見}
\bkmichishirube{2-7,4-6,8-8,力-1,力-2,力-3,力-15,電-1,電-2,電-4,電-6}{原-3,原-4,原-7,原-11,原-12,原-13}

\begin{goal}
\begin{itemize}[label={\color{bkTheme}●}, leftmargin=1.6\zw, itemsep=0pt, topsep=1pt]
\item \bklevel{1}\enspace 原子は、まん中の小さな原子核と、その周りの電子からできていることを、実験の結果から説明できる。
\item \bklevel{3}\enspace トムソンの実験で電子の比電荷を、ミリカンの実験で電気素量を、力のつり合いと円運動から求められる。
\item \bklevel{3}\enspace ラザフォードの実験から、原子核がとても小さく重いことを、いちばん近づく距離で見積もれる。
\item \bklevel{4}\enspace ねらいの距離と散乱角の関係を使って、大きく曲がる α 線の割合を計算できる。
\end{itemize}
\end{goal}

\phkoushiki{%
\phk{真空放電の陰極線は、電子の流れ}{中学理科}{ph:原-1:sec1}
\phk{原子は電子と原子核、原子核は陽子と中性子}{中学理科}{ph:原-1:sec6}
\phk{原子はとても小さい（約 0.1〜0.3\,nm）}{中学理科}{ph:原-1:sec7}
\phk{比電荷 $v=E/B$、$e/m=v/(rB)$}{高校物理}{ph:原-1:thomson}
\phk{電気素量 $qE=mg$（ミリカン）}{高校物理}{ph:原-1:millikan}
\phk{ラザフォードの原子模型}{高校物理}{ph:原-1:rutherford}
}

\begin{hatena}
\noindent\begin{minipage}[c]{0.56\linewidth}\raggedright
机も、水も、空気も、わたしたちの体も、原子という小さな粒が集まってできています。原子は、ぎっしりつまった小さな玉のようなものだと思っていませんか。

ところが、原子の中をのぞいてみると、まん中の「芯」はとても小さく、そのほかの所は、ほとんど何もないすき間なのだそうです。では、原子の中は、どのくらいからっぽなのでしょう。そして、見ることもできないほど小さな原子の中が、どうしてわかったのでしょう。
\end{minipage}\hfill
\begin{minipage}[c]{0.40\linewidth}\centering
\begin{tikzpicture}[x=1mm, y=1mm]
  \fill[bkYellow!20] (0,0) circle (21);
  \draw[bkBlue!60!black, line width=0.6pt, dash pattern=on 2pt off 1.5pt] (0,0) circle (21);
  \foreach \a in {20,95,170,250,320} {\txbel{(\a:14)}{1.6}}
  \fill[bkRed!75!black] (0,0) circle (0.45);
  \draw[bkInk, line width=0.4pt] (0.4,0.4) -- (8,7);
  \node[txb lab, anchor=south west, inner sep=1pt] at (8,7) {芯？};
  \node[txb lab, anchor=north] at (0,-22) {原子の中は、どうなっている？};
\end{tikzpicture}
\end{minipage}
\end{hatena}

\noindent 物理シリーズの『原子』は、この章から始まります。第1部（原-1〜4）の題は「飛び出す」です。電子・光・原子核が「飛び出して」くる実験から、原子の中のようすを読みとります。この章では、19世紀の終わりから20世紀のはじめにかけて、電子と原子核が見つかった実験をたどります。どの実験も、物理シリーズの『力学』『電磁気学』で学んだ、力のつり合い・円運動・エネルギーの保存だけで読みとけます。

\begin{yattemiyou}[10円玉と水1滴に、原子はいくつある？]
\textbf{用意するもの}：電卓、10円玉（重さは約\qty{4.5}{g}）、スポイト、はかりか計量スプーン\par\vspace{2pt}
\begin{enumerate}[label={\bkcircled[bkGreen]{\arabic*}}, leftmargin=2.4\zw, itemsep=1pt, topsep=0pt]
\item 10円玉は、ほとんど銅でできています。銅\qty{1}{mol}は\qty{63.55}{g}で、その中に原子が $6.02\times10^{23}$ 個（アボガドロ定数）あります。10円玉の重さ÷63.55×$6.02\times10^{23}$ を計算する。
\item スポイトで水を20滴落とし、約\qty{1}{g}になることを確かめる（計量スプーンの小さじ1杯は\qty{5}{g}）。水1滴（約\qty{0.05}{g}）の分子の数を、水\qty{1}{mol}＝\qty{18.0}{g}として計算する。
\item ②の分子を、1個\qty{0.31}{nm}（例題8）ずつ1列に並べたら、何 km になるか計算する。
\end{enumerate}
\end{yattemiyou}

\noindent やってみようの答え：①10円玉の原子は $4.5\div63.55\times6.02\times10^{23}\approx4.3\times10^{22}$ 個です。②水1滴の分子は $0.05\div18.0\times6.02\times10^{23}\approx1.7\times10^{21}$ 個。③1列に並べると、$1.7\times10^{21}\times0.31\times10^{-9}\,\mathrm{m}\approx5\times10^{11}$\,m、つまり約5億 km で、地球と太陽の間（約1.5億 km）の3倍以上です。水1滴の中にも、それほどたくさんの分子があるのです。1個1個の原子は、それほど小さく、目でも、光の顕微鏡でも見えません（物理シリーズの『波動』（波-13）で学んだように、光の顕微鏡で見分けられるのは約 \qty{0.2}{\micro m} まで）。この章の科学者たちは、見えない原子の中を、電気と磁場と、飛んでくる粒を使って調べました。

\begin{omoidasou}
\bkrecall{電-1}{点電荷どうしの間のクーロン力は $F=kq_1q_2/r^2$（$k=\qty{8.99e9}{N.m^2/C^2}$）。電気素量は $e=\qty{1.60e-19}{C}$。}
\bkrecall{電-2}{電荷 $q$ を電圧 $V$ で加速すると、運動エネルギーが $qV$ 増える。電子を\qty{1}{V}で加速したときのエネルギーが \qty{1}{eV}（$\qty{1.60e-19}{J}$）。}
\bkrecall{電-6}{磁場 $B$ に垂直に速さ $v$ で入った電荷 $q$ は、ローレンツ力 $qvB$ を受けて円を回る。半径 $r=mv/(qB)$（電-6 ⑩）。}
\bkrecall{力-1}{速さ $v$、半径 $r$ の等速円運動の加速度は、中心へ向かって $v^2/r$。}
\bkrecall{力-2}{小さな粒がゆっくり動くときの空気の抵抗は $f=bv$、$b=6\pi\eta a$（$\eta$ は空気の粘度、$a$ は粒の半径）。重力とつり合うと、一定の速さ（終端速度）で落ちる（力-2 例題12）。}
\bkrecall{力-3}{力が保存力だけなら、運動エネルギーと位置エネルギーの和は変わらない。}
\bkrecall{2-7}{$10^{-9}$ は10億分の1、$10^{-15}$ は1000兆分の1。$10^a\times10^b=10^{a+b}$。}
\end{omoidasou}

\noindent この章の記号：電場は $E$、運動エネルギーは $K$ と書きます（この章ではエネルギーを $E$ と書きません）。油滴の半径は $a$、空気の粘度は $\eta$（イータ）と書き、円運動の半径 $r$ と分けます。長さの単位は、ナノメートル（$\qty{1}{nm}=10^{-9}$\,m）と、フェムトメートル（$\qty{1}{fm}=10^{-15}$\,m。1 nm の100万分の1）を使います。

%% ======================== 物理の道具 ========================
\bksection{doguu}{物理の道具}

\bksubsection{1}{陰極線：真空の管の中を飛ぶもの}
\label{ph:原-1:sec1}
ガラスの管から空気をぬいていき、両はしの金属の電極に数千〜数万 V の電圧をかけると、管の中が光ったり、ガラスの壁が緑色に光ったりします。この現象を\textbf{真空放電}といいます。中学の理科では、この管（クルックス管）の中に、－極（\textbf{陰極}）から＋極（陽極）へ向かって何かが飛んでいることを、影や光の筋で見ました。この「何か」を\textbf{陰極線}といいます。\bkidx{いんきょくせん}{陰極線}

\begin{bkfigure}
\begin{tikzpicture}[x=1mm, y=1mm]
  \draw[bkInk, line width=0.8pt, fill=bkBlue!5] (0,0) rectangle (100,24);
  \fill[bkGray!60] (2,5) rectangle (5,19); \node[txb lab, anchor=south] at (3.5,25) {陰極（－）};
  \fill[bkGray!60] (20,7) rectangle (22,10.5) (20,13.5) rectangle (22,17); \node[txb lab, anchor=south] at (21,25) {陽極（＋）};
  \fill[txb plate] (35,18) rectangle (60,19.5); \fill[txb plate] (35,4.5) rectangle (60,6);
  \node[txb lab, anchor=south] at (47.5,25) {上の板（＋）};
  \node[txb lab, anchor=north] at (47.5,-1) {下の板（－）};
  \draw[txb beam] (5,12) -- (35,12) .. controls (50,12) and (55,13.5) .. (60,15) -- (99,20.4);
  \fill[bkGreen!60] (99,4) rectangle (100,22);
  \node[txb lab, anchor=west, align=left] at (102,20) {ガラスが\\光る点};
  \node[txb lab, anchor=west, align=left] at (102,8) {板がないとき\\の点};
  \draw[bkOrange!85!black, line width=0.8pt, dash pattern=on 3pt off 1.5pt] (60,12) -- (99,12);
  \draw[txb F] (48,9) -- (48,15.5);
  \node[txb lab, anchor=west] at (49,9.5) {力};
\end{tikzpicture}
\bkfigcaption{陰極線の管（模式図）。陽極の穴を通りぬけた陰極線は、上下の板に電圧をかけると、\textbf{＋の板のほうへ}曲がる（実線）。電圧がないとまっすぐ進む（破線）。}\label{ph:原-1:fig-inkyoku}
\end{bkfigure}

\noindent 陰極線は、何でできているのでしょう。19世紀の後半、光のような波だと考える人と、小さな粒の流れだと考える人がいました。手がかりは、次の3つの実験です。
\begin{enumerate}[label={\bkcircled[bkBlue]{\arabic*}}, leftmargin=2.4\zw, itemsep=1pt, topsep=2pt]
\item 陰極線の通り道に金属の板を置くと、ガラスに板の影ができる。陰極線は、陰極からまっすぐ進む。
\item 図\ref{ph:原-1:fig-inkyoku}のように、上下の板に電圧をかけると、陰極線は\textbf{＋の板のほうへ}曲がる。＋に引かれるのは、負の電荷をもつものです（電-1）。
\item 磁石を近づけると曲がる。曲がる向きは、負の電荷が陰極から陽極へ動いているとして、ローレンツ力（電-6）で考えた向きと同じ。
\end{enumerate}
②と③から、陰極線は\textbf{負の電荷をもつ粒の流れ}だとわかります。陰極の金属の種類を変えても、管に入れる気体の種類を変えても、同じ性質の粒が出てきます。この粒が\textbf{電子}です。\bkidx{でんし}{電子}金属の導線の中を流れる電流も、この電子の流れです（電-4 ①）。電流の向きは、電子が動く向きと反対だと決めたのでした。

\begin{phmijika}[ブラウン管のテレビと、蛍光灯]
昔のテレビ（ブラウン管）は、陰極線の管そのものです。管の奥の電子銃から電子を飛ばし、磁場で向きを変えて、画面の内側にぬった蛍光体の点に当てて光らせていました。1秒に何十枚もの絵を、電子の筋でなぞって描いていたのです。蛍光灯やネオンサインも、管の中の気体の中を電子が走り、気体の原子にぶつかって光を出させるしくみです（光る理由は原-4）。
\end{phmijika}

\bksubsection{2}{トムソンの実験：電子の比電荷}
\label{ph:原-1:sec2}
陰極線が負の電荷の粒だとすると、次に知りたいのは、粒1個の電荷 $q$ と質量 $m$ です。ところが、電場や磁場の中での粒の動き方は、運動方程式
\[ m\vec a=q\vec E+q\vec v\times\vec B \]
で決まります。両辺を $m$ でわると、加速度は $q/m$ だけで決まります。そのため、曲がり方をはかっても、$q$ と $m$ を別々に知ることはできず、わかるのは比 $q/m$ だけです。この比を\textbf{比電荷}といいます。\bkidx{ひでんか}{比電荷}

1897年、イギリスのトムソンは、陰極線に電場と磁場を同時にかける方法で、比電荷をはかりました（1906年にノーベル物理学賞）。図\ref{ph:原-1:fig-thomson}のように、上下の板の間に下向きの電場 $E$ をつくり、同じ場所に、紙のおくへ向かう磁場 $B$ をかけます。

\begin{bkfigure}
\begin{tikzpicture}[x=1mm, y=1mm]
  \fill[txb plate] (10,40) rectangle (60,42); \fill[txb plate] (10,0) rectangle (60,2);
  \node[txb lab, anchor=south] at (35,42.5) {＋の板};
  \node[txb lab, anchor=north] at (35,-0.5) {－の板};
  \foreach \x in {14,24,44,54} {\foreach \y in {7,35} {\txbin{(\x,\y)}{1.6}}}
  \node[txb lab, anchor=west] at (61,6) {磁場 $B$（紙のおくへ）};
  \draw[bkGreen!50!black, line width=0.7pt, -{Stealth[length=2mm]}] (66,38) -- (66,27);
  \node[txb lab, anchor=west] at (67,33) {電場 $E$（下向き）};
  \draw[txb beam] (-8,21) -- (90,21);
  \fill[bkBlue!70!black] (-8,21) circle (0.8);
  \node[txb lab, anchor=south] at (-4,22) {電子};
  \txbel{(34,21)}{1.6}
  \draw[txb F] (34,23) -- (34,32);
  \node[txb lab, anchor=west] at (35.5,29) {電場の力 $eE$};
  \draw[txb Fb] (34,19) -- (34,10);
  \node[txb lab, anchor=west] at (35.5,13) {磁場の力 $evB$};
  \draw[bkInk, line width=0.7pt, -{Stealth[length=2mm]}] (19,23) -- (27,23);
  \node[txb lab, anchor=south] at (23,23.5) {$v$};
\end{tikzpicture}
\bkfigcaption{トムソンの実験の考え方。電子は負なので、電場の力は電場と反対の上向き（実線の矢印）、磁場の力 $e\vec v\times\vec B$ の向きは、負の電荷なので反対になって下向き（破線の矢印）。2つの力がつり合う速さの電子だけが、まっすぐ通りぬける。}\label{ph:原-1:fig-thomson}
\end{bkfigure}

\noindent 電子の電荷の大きさを $e$ と書くと、電子が受ける力は、電場から上向きに $eE$、磁場から下向きに $evB$ です（電-6 ⑩）。2つの力がちょうどつり合うと、電子はまっすぐ進みます。
\pheq{eE=evB\qquad\text{より}\qquad v=\frac{E}{B}}{m/s ＝ (V/m)÷T}
\label{ph:原-1:thomson}$e$ が両辺で消えるので、電子の電荷や質量を知らなくても、電場と磁場の強さだけで\textbf{速さがわかる}のです。次に、電場を切って磁場だけにすると、電子は円をえがきます。円運動の式（力-1）に、中心へ向かう力 $evB$ を入れると、
\[ m\,\frac{v^2}{r}=evB\qquad\text{より}\qquad r=\frac{mv}{eB} \]
です。円の半径 $r$ をはかれば、比電荷が求まります。
\pheq{\frac{e}{m}=\frac{v}{rB}=\frac{E}{rB^2}}{C/kg ＝ (m/s)÷(m×T)}

\begin{phyomi}[$\dfrac{e}{m}=\dfrac{v}{rB}$]
「速く飛ぶほど、曲がりにくい（$r$ が大きい）。電荷が大きいほど、磁場の力が大きく、よく曲がる。重いほど、曲がりにくい」。この3つをまとめると、半径は $mv/(eB)$ になります。半径・速さ・磁場の3つをはかれば、$e/m$ が決まります。
\end{phyomi}

\noindent トムソンが求めた陰極線の比電荷は、水素のイオン（水素の原子から電子が1個とれたもの）の比電荷の、およそ1000倍以上でした。電荷が同じくらいなら、陰極線の粒は、いちばん軽い水素の原子より、ずっと軽いことになります。原子より小さい粒が見つかったのです。今の値は
\pheq{\frac{e}{m_{\mathrm e}}=\qty{1.76e11}{C/kg}}{C/kg}
で、水素のイオン（陽子）の比電荷 $\qty{9.58e7}{C/kg}$ の約1840倍です。\qty{1}{kg}の電子は、$1.76\times10^{11}$\,C という、とてつもない電荷をもつことになります。

\begin{phbaisuru}{磁場 $B$}
同じ速さの電子なら、半径 $r=mv/(eB)$ は半分になります。磁場を強くするほど、小さな円を回ります。下のグラフは、$v=\qty{2.0e7}{m/s}$ の電子の半径です（例題2の速さ）。
\par\smallskip\noindent\begin{tikzpicture}
\begin{axis}[bkplot, width=96mm, height=40mm, xmin=0, xmax=4.2, ymin=0, ymax=0.25, xtick={0.5,1,2,3,4}, ytick={0.05,0.10,0.15,0.20}, yticklabel style={/pgf/number format/fixed}, xlabel={磁場 $B$（mT）}, ylabel={半径 $r$（m）}]
  \addplot[bkcurve, domain=0.46:4.1, samples=80] {0.1137/x};
  \fill[bkRed] (axis cs:1,0.1137) circle (1.8pt) (axis cs:2,0.0569) circle (1.8pt);
  \node[font=\sffamily\footnotesize, anchor=south west, text=bkRed!75!black] at (axis cs:1.1,0.12) {\qty{1}{mT}で\qty{11}{cm}};
  \node[font=\sffamily\footnotesize, anchor=south west, text=bkRed!75!black] at (axis cs:2.1,0.06) {2倍で半分};
\end{axis}
\end{tikzpicture}
\end{phbaisuru}

\begin{phflow}[ローレンツ力から比電荷へ]
\phflowrow{$m\vec a=q(\vec E+\vec v\times\vec B)$}{電場と磁場が垂直\\まっすぐ進む}{$v=\dfrac{E}{B}$}
\phflowrow{$m\vec a=q\vec v\times\vec B$}{電場を切る\\等速円運動}{$r=\dfrac{mv}{eB}$}
\phflowrow{$v=\dfrac{E}{B}$、$r=\dfrac{mv}{eB}$}{2つの式から\\$v$ を消す}{$\dfrac{e}{m}=\dfrac{E}{rB^2}$}
\end{phflow}

\begin{phmodel}[トムソンの実験の電子]
\phignore{電子にはたらく重力（$10^{-29}$\,N の桁で、電場の力 $10^{-15}$\,N の桁よりずっと小さい）。管に残った気体の分子とのぶつかり。}
\phassume{電場と磁場は、板の間で一様。電子はみな同じ速さ。光の速さよりずっと遅い。}
\phlimit{電子が光の速さに近いとき（数十 kV 以上で加速したとき。原-11）。板のはしで電場が弱まる所。}
\end{phmodel}

\bksubsection{3}{ミリカンの油滴：電気素量}
\label{ph:原-1:sec3}
比電荷がわかっても、電子1個の電荷 $e$ と質量 $m$ は、別々にはわかりません。電荷そのものをはかったのが、アメリカのミリカンです（1909〜1913年。1923年にノーベル物理学賞）。

図\ref{ph:原-1:fig-millikan}のように、間かく $d_{\mathrm p}$ の2枚の平らな板に電圧 $V$ をかけると、板の間には一様な電場 $E=V/d_{\mathrm p}$ ができます（電-2 ⑥）。負の電荷 $-q$ をもつ油滴は、電場と反対向きの上向きに力 $qE$ を受けます。電圧を調整して、油滴が止まって浮かんでいるようにすると、上向きの電場の力と下向きの重さがつり合っています。
\pheq{qE=mg\qquad\text{より}\qquad q=\frac{mg}{E}=\frac{mgd_{\mathrm p}}{V}}{C ＝ N÷(V/m)}
\label{ph:原-1:millikan}油滴の質量 $m$ がわかれば、電荷 $q$ が求まります（$m$ の求め方は④）。

\begin{bkfigure}
\begin{tikzpicture}[x=1mm, y=1mm]
  \fill[txb plate] (0,32) rectangle (60,34.5); \fill[txb plate] (0,0) rectangle (60,2.5);
  \fill[white] (27,32) rectangle (33,34.5);
  \node[txb lab, anchor=south west] at (0,34.8) {上の板（＋）};
  \node[txb lab, anchor=north west] at (0,-0.4) {下の板（－）};
  \draw[bkInk, line width=0.5pt, fill=bkGray!20] (24,45) -- (36,45) -- (33,38) -- (27,38) -- cycle;
  \node[txb lab, anchor=west] at (37,43) {霧ふき（油の霧を入れる）};
  \foreach \p in {(29,36.5),(31.5,37.5),(30,39.5)} {\fill[bkOrange!80!black] \p circle (0.5);}
  \filldraw[fill=bkOrange!70, draw=bkInk, line width=0.4pt] (30,17) circle (2.2);
  \draw[txb F] (30,19.5) -- (30,28.5);
  \node[txb lab, anchor=west] at (31.5,26) {電場の力 $qE$};
  \draw[txb Fb] (30,14.5) -- (30,5.5);
  \node[txb lab, anchor=west] at (31.5,7.5) {重さ $mg$};
  \draw[txb dim] (-4,2.5) -- (-4,32);
  \node[txb lab, anchor=east] at (-5,17) {$d_{\mathrm p}$};
  \draw[bkInk, line width=0.6pt] (82,17) circle (4); \draw[bkInk, line width=0.6pt] (78,17) -- (64,17);
  \node[txb lab, anchor=north] at (78,12) {顕微鏡でのぞく};
  \draw[bkInk, line width=0.6pt] (60,33.2) -- (100,33.2) -- (100,26);
  \draw[bkInk, line width=0.6pt] (60,1.2) -- (100,1.2) -- (100,8);
  \draw[bkInk, line width=0.6pt, fill=white] (95,8) rectangle (105,26);
  \node[txb lab, align=center] at (100,17) {電\\圧\\$V$};
\end{tikzpicture}
\bkfigcaption{ミリカンの実験（模式図）。霧ふきで細かい油の粒（油滴）を板の間に入れ、顕微鏡で1つの油滴を見ながら電圧を調整する。油滴は、霧ふきの口を通るときにこすれて、電荷をもつ。}\label{ph:原-1:fig-millikan}
\end{bkfigure}

ミリカンは、何千もの油滴の電荷をはかりました。すると、どの油滴の電荷も、ある数の\textbf{整数倍}になっていました。$1.5$ 倍や $2.3$ 倍の電荷をもつ油滴は、1つもなかったのです。この最小の単位が\textbf{電気素量} $e$ です。\bkidx{でんきそりょう}{電気素量}電荷は、水のように好きな量に分けられるのではなく、$e$ ずつのとびとびの量でしかやりとりできないのです。今では、$e$ はちょうど
\[ e=1.602\,176\,634\times10^{-19}\ \mathrm{C} \]
と決められています（2019年から、この値を使って電流の単位アンペアを決めている：電-6）。この本では $e=\qty{1.60e-19}{C}$ を使います。

\begin{bkfigure}
\begin{tikzpicture}[x=1mm, y=1mm]
  \draw[bkInk, line width=0.7pt, -{Stealth[length=2mm]}] (0,0) -- (120,0);
  \foreach \k in {0,...,7} {\draw[bkInk, line width=0.5pt] ({\k*16},-1.5) -- ({\k*16},1.5); \node[txb lab, anchor=north] at ({\k*16},-2) {\k};}
  \node[txb lab, anchor=north east] at (122,-6.5) {目もりの単位：$1.60\times10^{-19}$\,C};
  \foreach \k/\c in {1/2,2/4,3/3,4/2,5/1,6/1,7/1} {\foreach \j in {1,...,\c} {\fill[bkBlue!70!black] ({\k*16},{\j*3.4}) circle (1.3);}}
  \node[txb lab, anchor=west] at (0,24) {はかった油滴の電荷（1つの点が1つの油滴）};
  \draw[bkRed!80!black, line width=0.6pt, dash pattern=on 2pt off 1.5pt] (24,1.5) -- (24,14) (56,1.5) -- (56,14);
  \node[txb lab, text=bkRed!80!black, anchor=south] at (24,14.5) {1.5倍};
  \node[txb lab, text=bkRed!80!black, anchor=south] at (56,14.5) {3.5倍};
\end{tikzpicture}
\bkfigcaption{油滴の電荷は、$e$ の整数倍の所にだけ集まる（点の高さは、同じ電荷の油滴の数）。1.5倍や3.5倍の所には1つもない。}\label{ph:原-1:fig-seisu}
\end{bkfigure}

\begin{phyomi}[$qE=mg$]
「油滴を、電気の力で持ち上げて、重さとつり合わせる」。手のひらに乗せた物の重さを、手で押し上げる力ではかるのと同じです。電荷が2倍の油滴は、半分の電圧で浮かびます。
\end{phyomi}

\begin{phbaisuru}{油滴の半径 $a$}
油滴の質量は、体積 $\tfrac43\pi a^3$ に比例するので8倍になります。同じ電荷なら、浮かべるのに要る電圧も8倍です。ミリカンが半径 \qty{1}{\micro m} ほどの小さな油滴を選んだのは、電荷が $e$ の数倍しかないときに、ちょうどよい電圧（数百 V）で浮かぶからです（例題4）。
\end{phbaisuru}

\bksubsection{4}{油滴の大きさ：終端速度とストークスの抵抗}
\label{ph:原-1:sec4}
③の式で電荷を求めるには、油滴の質量 $m$ が要ります。でも、半径 \qty{1}{\micro m} ほどの油滴を、はかりにのせることはできません。ミリカンは、油滴が\textbf{落ちる速さ}から半径を求めました。

電圧を切ると、油滴は落ち始めます。小さな粒がゆっくり動くときの空気の抵抗は、速さに比例します（物理シリーズの『力学』（力-2 例題12）の霧の粒と同じ。$6\pi$ の出どころは力-15）。
\[ f=6\pi\eta av \]
$\eta$（イータ）は空気の粘度（20\,℃で $\qty{1.81e-5}{Pa.s}$）、$a$ は油滴の半径です。落ちる速さが増えると抵抗も増え、やがて重さとつり合って、一定の速さ $v_{\mathrm t}$（\textbf{終端速度}）で落ちるようになります。\bkidx{しゅうたんそくど}{終端速度}油滴の質量は、密度 $\rho$ と体積 $\tfrac43\pi a^3$ で $m=\tfrac43\pi a^3\rho$ なので、
\[ \frac43\pi a^3\rho g=6\pi\eta av_{\mathrm t} \qquad\text{より}\qquad a=\sqrt{\frac{9\eta v_{\mathrm t}}{2\rho g}} \]
\label{ph:原-1:stokes}です。顕微鏡の目もりで、油滴が決まった長さを落ちる時間をストップウォッチではかれば $v_{\mathrm t}$ がわかり、半径 $a$、質量 $m$ の順に求まります。半径 \qty{1}{\micro m} の油滴の終端速度は、1秒に約 \qty{0.1}{mm} です（例題6）。電圧をかけて止めるときの重さと、電圧を切って落とすときの速さの、2つの測定を組み合わせるのが、ミリカンの方法のかなめです。

\begin{phbaisuru}{落ちる油滴の半径 $a$}
重さは $a^3$ に比例して8倍、空気の抵抗は $a$ に比例するので、つり合う速さ（終端速度 $v_{\mathrm t}=2\rho ga^2/(9\eta)$）は $a^2$ に比例して4倍になります。大きな油滴ほど速く落ちます。雨つぶが霧の粒より速く落ちるのと同じです（ただし、大きな雨つぶでは抵抗の式が変わる：力-15）。
\par\smallskip\noindent\begin{tikzpicture}
\begin{axis}[bkplot, width=88mm, height=36mm, xmin=0, xmax=2.2, ymin=0, ymax=0.48, xtick={0.5,1,1.5,2}, ytick={0.1,0.2,0.3,0.4}, yticklabel style={/pgf/number format/fixed}, xlabel={半径 $a$（\unit{\micro m}）}, ylabel={$v_{\mathrm t}$（mm/s）}]
  \addplot[bkcurve, domain=0:2.1, samples=60] {0.1047*x^2};
  \fill[bkRed] (axis cs:1,0.1047) circle (1.8pt) (axis cs:2,0.4188) circle (1.8pt);
  \node[font=\sffamily\footnotesize, anchor=south east, text=bkRed!75!black] at (axis cs:0.95,0.12) {\qty{1}{\micro m}で 0.10};
  \node[font=\sffamily\footnotesize, anchor=east, text=bkRed!75!black] at (axis cs:1.9,0.40) {2倍で4倍};
\end{axis}
\end{tikzpicture}
\end{phbaisuru}

\begin{phmodel}[ミリカンの油滴]
\phignore{空気の浮力（油の密度の約700分の1）。油滴が蒸発して小さくなること。}
\phassume{油滴は球。空気の抵抗はストークスの抵抗 $6\pi\eta av$。落ち始めてすぐに終端速度になる（$10^{-5}$ 秒ほどで近づく）。}
\phlimit{油滴が空気の分子の間かく（約 \qty{0.07}{\micro m}）に近いほど小さいとき（抵抗が小さくなる。ミリカンはこの直しも入れた）。空気の分子にぶつかってふらふら動く（ブラウン運動）のが目立つとき。}
\end{phmodel}

\bksubsection{5}{ラザフォードの散乱：原子核の発見}
\label{ph:原-1:sec5}
電子は、どの原子にもふくまれています。でも、原子全体は電気を帯びていません。だから、原子の中には、電子の負の電荷を打ち消すだけの\textbf{正の電荷}があるはずです。また、電子はとても軽いので、原子の重さのほとんどは、この正の部分がもっているはずです。では、正の電荷は、原子の中にどのように広がっているのでしょう。

20世紀のはじめには、「正の電荷が原子の大きさの球全体に一様に広がり、その中に電子がちらばっている」という考え（トムソンの模型）が有力でした。これを確かめるために、イギリスのマンチェスターで、ラザフォードの指導のもと、ガイガーとマースデンは、ラジウムなどが出す \textbf{α 線}（アルファ線）を、うすい金はくに当てる実験をしました（1909年）。α 線は、正の電荷 $2e$ をもち、電子の約7300倍の重さの粒（ヘリウムの原子核：原-12）が、とても速く飛んでいるものです。


図\ref{ph:原-1:fig-rutherford}のような装置で数えた結果は、おどろくべきものでした。ほとんどの α 線は、金はくをまっすぐ通りぬけました。ところが、ごくまれに、$90^\circ$ より大きく曲がって、来た向きへはね返されるものがあったのです。トムソンの模型のように、正の電荷が原子全体にうすく広がっていると、電荷の広がった球の中で受ける力は弱く、重くて速い α 線を大きく曲げることはできません。1911年、ラザフォードは、次のように考えました（1908年にノーベル化学賞を受けたのは、これより前の放射能の研究）。

\begin{bkcenter}
\begin{minipage}{0.92\linewidth}
原子の正の電荷と、質量のほとんどは、原子のまん中の\textbf{とても小さな所}に集まっている。これを\textbf{原子核}という。電子は、その周りの広い空間にいる。
\end{minipage}
\end{bkcenter}

\begin{bkfigure}
\begin{tikzpicture}[x=1mm, y=1mm]
  \draw[bkInk, line width=0.5pt, fill=bkGray!60] (0,10) rectangle (12,26); \fill[white] (8,16) rectangle (12,20);
  \fill[bkRed!75!black] (4,18) circle (1.2);
  \node[txb lab, anchor=north, align=center] at (6,9) {α 線を出す物\\（鉛の箱の中）};
  \draw[txb beam] (12,18) -- (44,18);
  \fill[bkYellow!70!black] (44,6) rectangle (45.2,30);
  \node[txb lab, anchor=south] at (44.6,31) {金はく};
  \draw[bkGreen!40!black, line width=1pt] (44.6,18) ++(-50:26) arc (-50:210:26);
  \node[txb lab, anchor=west, align=left] at (74,4) {光る幕（α 線が当たると\\小さく光る）を回して数える};
  \draw[txb beam] (45.2,18) -- (70.6,18);
  \draw[txb beam] (45.2,18.2) -- (70.0,23.4);
  \draw[txb beam] (45.2,17.8) -- (67.8,6.2);
  \draw[bkRed!80!black, line width=1.3pt, -{Stealth[length=2.4mm]}] (44,18.4) -- (25.5,29);
  \node[txb lab, text=bkRed!80!black, anchor=east, align=right] at (17,36) {ごくまれに\\大きくはね返る};
  \node[txb lab, anchor=west, align=left] at (74,28) {ほとんどは\\まっすぐ};
\end{tikzpicture}
\bkfigcaption{ガイガーとマースデンの実験（模式図）。真空の容器の中で行う。α 線が当たった光る幕の点を、暗い部屋で顕微鏡をのぞいて1つずつ数えた。}\label{ph:原-1:fig-rutherford}
\end{bkfigure}

\noindent\label{ph:原-1:rutherford}原子核が小さければ、原子核のすぐそばを通る α 線は、近いので大きなクーロン力を受けて、大きく曲げられます。原子核はとても小さいので、そのすぐそばを通る α 線は、ほんの少ししかありません。これで、「ほとんどはまっすぐ、ごくまれに大きくはね返る」が説明できます。\bkidx{げんしかく}{原子核}

\Needspace{6\baselineskip}\phboxhead{原子核の大きさを見積もる}
原子核の大きさは、どのくらいでしょう。α 線が、金の原子核に\textbf{正面から}向かう場合を考えます。α 線（電荷 $2e$）と金の原子核（電荷 $Ze$、金は $Z=79$）は、どちらも正なので退け合い、α 線は近づくほど遅くなって、いちばん近づいた所で一瞬止まり、はね返されます。金の原子核は α 線の約50倍重いので、動かないとします。はじめの運動エネルギー $K$ が、いちばん近づいた距離 $d$ でのクーロン力の位置エネルギー（電-2 ⑤）に全部変わるので、エネルギーの保存から
\pheq{K=k\,\frac{(2e)(Ze)}{d}\qquad\text{より}\qquad d=\frac{2Zke^2}{K}}{m ＝ (N·m$^2$/C$^2$)×C$^2$÷J}
です。原子の世界では、エネルギーを MeV（100万 eV）、長さを fm で書くと便利です。
\[ ke^2=\frac{8.99\times10^9\times(1.60\times10^{-19})^2}{1.60\times10^{-13}\ \mathrm{J/MeV}}\ \mathrm{J\cdot m}\approx1.44\times10^{-15}\ \mathrm{MeV\cdot m}=\qty{1.44}{MeV.fm} \]
ガイガーとマースデンの α 線は約 \qty{5}{MeV}〜\qty{8}{MeV} でした。$K=\qty{5.0}{MeV}$ なら
\[ d=\frac{2\times79\times1.44}{5.0}\approx\qty{45}{fm} \]
です（例題9）。α 線はここまで近づいてもクーロン力の式どおりに曲がったので、金の原子核の半径は \qty{45}{fm} より小さいことがわかります。原子の半径（約 \qty{0.14}{nm}＝\num{140000}\,fm）の3000分の1よりも小さいのです。

\begin{phyomi}[$d=\dfrac{2Zke^2}{K}$]
「速い（エネルギーの大きい）α 線ほど、原子核の近くまで入りこめる」。坂をボールが転がり上がるのと同じで、勢いが大きいほど、高い所（電気の位置エネルギーが大きい所＝原子核のすぐそば）まで上れます。原子核の電荷 $Ze$ が大きいほど、手前で止められます。
\end{phyomi}

\begin{phbaisuru}{α 線のエネルギー $K$}
いちばん近づく距離は半分になります。$d$ は $K$ に反比例します。エネルギーを上げていくと、ついに α 線が原子核にふれ、クーロン力だけでは説明できない曲がり方が始まります。そのエネルギーから、原子核の大きさをはかることができます（例題9のじぶんでとこう）。グラフの破線は、金の原子核の半径（約\qty{7}{fm}）です。
\par\smallskip\noindent\begin{tikzpicture}
\begin{axis}[bkplot, width=90mm, height=40mm, xmin=0, xmax=21, ymin=0, ymax=120, xtick={5,10,15,20}, ytick={20,45,80}, xlabel={α 線のエネルギー $K$（MeV）}, ylabel={$d$（fm）}]
  \addplot[bkcurve, domain=1.9:20.5, samples=80] {227.5/x};
  \fill[bkRed] (axis cs:5,45.5) circle (1.8pt) (axis cs:10,22.75) circle (1.8pt);
  \node[font=\sffamily\footnotesize, anchor=south west, text=bkRed!75!black] at (axis cs:5.5,47) {\qty{5}{MeV}で\qty{45}{fm}};
  \node[font=\sffamily\footnotesize, anchor=south west, text=bkRed!75!black] at (axis cs:10.5,27) {2倍で半分};
  \addplot[bkthin, dashed, domain=0:20.5] {7};
\end{axis}
\end{tikzpicture}
\end{phbaisuru}

\Needspace{5\baselineskip}\phboxhead{\bkhatenaref の答え}
\label{ph:原-1:hatena}今では、原子核の半径は、多くの原子核の測定から、質量数（陽子と中性子の数の和：⑥）を $A$ として、およそ $1.2\times A^{1/3}$\,fm だとわかっています（原-11）。金（$A=197$）では約 \qty{7.0}{fm} です。一方、金の原子の半径（となりの原子との中心の間の半分）は約 \qty{0.144}{nm} です。長さの比は
\[ \frac{7.0\times10^{-15}}{0.144\times10^{-9}}\approx4.8\times10^{-5}\qquad\text{（約2万分の1）} \]
体積の比は、その3乗で約 $1.1\times10^{-13}$（約9兆分の1）です（例題1）。原子を直径 \qty{100}{m} の競技場の大きさにすると、原子核は直径 約 \qty{5}{mm}、米つぶほどです。電子は、点のように小さいと考えられています。原子の体積のほとんどは、たしかに「からっぽ」なのです。それでも、原子の質量の 99.98\,\% は、この米つぶ（原子核）にあります。

\begin{bkfigure}
\begin{tikzpicture}[x=1mm, y=1mm]
  \fill[bkGreen!15] (0,0) ellipse (40 and 22);
  \draw[bkGreen!50!black, line width=0.8pt] (0,0) ellipse (40 and 22);
  \draw[bkGreen!50!black, line width=0.5pt] (0,0) ellipse (30 and 14);
  \fill[bkRed!75!black] (0,0) circle (0.5);
  \node[txb lab, anchor=north, align=center] at (0,-1.5) {原子核：米つぶ\\（直径 約\qty{5}{mm}）};
  \draw[txb dim] (-40,-26) -- (40,-26);
  \node[txb lab, anchor=north] at (0,-27) {原子：競技場（直径 \qty{100}{m}）};
  \node[txb lab, anchor=west, text width=40mm, align=left] at (46,6) {電子は、この競技場の客席の辺りを飛び回る、もっと小さな点。\par\vspace{2pt}競技場のほかの所は、からっぽ。};
\end{tikzpicture}
\bkfigcaption{原子を競技場の大きさにすると（金の原子。原子核は見やすいよう少し大きく描いた）。}\label{ph:原-1:fig-kyougijou}
\end{bkfigure}

\bksubsection{6}{原子核は陽子と中性子でできている}
\label{ph:原-1:sec6}
いちばん軽い原子である水素の原子核は、電荷 $+e$ をもつ1個の粒で、\textbf{陽子}とよばれます。\bkidx{ようし}{陽子}陽子は電子の約1840倍の重さです（②）。原子核の電荷は、陽子の数で決まり、陽子の数 $Z$ を\textbf{原子番号}といいます。原子全体は電気を帯びていないので、原子の中の電子の数も $Z$ 個です。原子番号が、元素の種類を決めます（水素1、ヘリウム2、炭素6、金79）。

ところが、原子核の重さをはかると、陽子 $Z$ 個の重さよりずっと重いのです。たとえばヘリウムの原子核は、電荷は $+2e$（陽子2個分）なのに、重さは陽子の約4倍あります。この「電荷のない重さ」の正体が、1932年にイギリスのチャドウィックが見つけた\textbf{中性子}です（1935年にノーベル物理学賞。見つけ方は原-11 ③）。\bkidx{ちゅうせいし}{中性子}中性子は電荷をもたず、重さは陽子とほぼ同じです。

\begin{bkfigure}
\begin{tikzpicture}[x=1mm, y=1mm]
  \foreach \p in {(1.8,1.8),(-1.8,-1.8)} {\filldraw[fill=bkGray!25, draw=bkInk, line width=0.4pt] \p circle (2.0);}
  \foreach \p in {(-1.8,1.8),(1.8,-1.8)} {\filldraw[fill=bkRed!70, draw=bkInk, line width=0.4pt] \p circle (2.0); \node[font=\sffamily\bfseries\scriptsize] at \p {＋};}
  \node[txb lab, anchor=north, align=center] at (0,-5) {ヘリウム4\\陽子2・中性子2};
  \begin{scope}[shift={(36,0)}]
  \foreach \x/\y in {-5.4/3.6,-1.8/3.6,1.8/3.6,5.4/3.6,-3.6/0,0/0,3.6/0,7.2/0,-5.4/-3.6,-1.8/-3.6,1.8/-3.6,5.4/-3.6} {\filldraw[fill=bkGray!25, draw=bkInk, line width=0.4pt] (\x,\y) circle (2.0);}
  \foreach \x/\y in {-5.4/3.6,1.8/3.6,0/0,7.2/0,-1.8/-3.6,5.4/-3.6} {\filldraw[fill=bkRed!70, draw=bkInk, line width=0.4pt] (\x,\y) circle (2.0); \node[font=\sffamily\bfseries\scriptsize] at (\x,\y) {＋};}
  \node[txb lab, anchor=north, align=center] at (0.9,-7) {炭素12\\陽子6・中性子6};
  \end{scope}
  \node[txb lab, anchor=west, text width=56mm, align=left] at (58,1) {＋の丸が陽子、無地の丸が中性子。\par\vspace{2pt}陽子と中性子をまとめて\textbf{核子}という。原子核の中の並び方は、絵のように決まってはいない。};
\end{tikzpicture}
\bkfigcaption{原子核のつくり（模式図）。陽子と中性子の大きさは約 \qty{1}{fm}。}\label{ph:原-1:fig-kakushi}
\end{bkfigure}

\noindent 陽子の数 $Z$ と中性子の数 $N$ の和 $A=Z+N$ を\textbf{質量数}といいます。原子の質量は、ほぼ質量数に比例します。同じ元素（同じ $Z$）でも、中性子の数がちがう原子があり、これを\textbf{同位体}といいます。\bkidx{どういたい}{同位体}炭素には、中性子が6個の炭素12（$^{12}$C）、7個の炭素13（$^{13}$C）、8個の炭素14（$^{14}$C）があります。同位体は、電子の数が同じなので、化学的な性質はほとんど同じです。でも重さがちがうので、⑪の質量分析で分けられます。原子核の中で、正の陽子どうしが退け合うのに、なぜばらばらにならないのか（核力）、炭素14 のような原子核がなぜこわれるのか（放射能）は、第4部（原-11・原-12）で学びます。

\bksubsection{7}{原子の大きさ：アボガドロ定数}
\label{ph:原-1:sec7}
原子1個の大きさは、物質の密度と、\qty{1}{mol}にふくまれる粒の数（\textbf{アボガドロ定数} $N_{\mathrm A}=6.02\times10^{23}$\,/mol）からわかります。\bkidx{あぼがどろていすう}{アボガドロ定数}たとえば水は、\qty{1}{mol}が\qty{18.0}{g}、\qty{1}{cm^3}が\qty{1.00}{g}なので、\qty{1}{mol}の体積は\qty{18.0}{cm^3}です。そこに $6.02\times10^{23}$ 個の分子があるので、分子1個が占める場所の体積と、それを立方体とみたときの1辺は、
\[ \frac{18.0\times10^{-6}\ \mathrm{m^3}}{6.02\times10^{23}}\approx3.0\times10^{-29}\ \mathrm{m^3},\qquad \sqrt[3]{3.0\times10^{-29}}\ \mathrm{m}\approx\qty{0.31}{nm} \]
です（例題8）。原子や小さな分子の大きさは、どれも \qtyrange{0.1}{0.3}{nm} くらいです。

では、$N_{\mathrm A}$ そのものは、どうやってはかるのでしょう。1つの方法が、電気分解です。銀のイオン Ag$^+$ を電気分解すると、電子1個で銀の原子1個が出てきます。そのため、銀\qty{1}{mol}を出すのに流れる電荷（\textbf{ファラデー定数}、\qty{96485}{C/mol}：電-4）は、電子 $N_{\mathrm A}$ 個分の電荷です。
\pheq{\num{96485}\ \mathrm{C/mol}=N_{\mathrm A}e\qquad\text{より}\qquad N_{\mathrm A}=\frac{\num{96485}}{1.60\times10^{-19}}\approx6.02\times10^{23}\ \mathrm{/mol}}{1/mol ＝ (C/mol)÷C}
ファラデー定数は、1830年代から電気分解の重さではかられていました。ミリカンが $e$ をはかったことで、$N_{\mathrm A}$ がはじめて正確に決まり、原子1個の大きさや重さが、はっきりした数で言えるようになったのです（ミリカンの1913年の論文の題も「電気素量とアボガドロ定数について」）。

\bksubsection{8}{散乱の式：ねらいの距離と曲がる角}
\label{ph:原-1:sec8}
⑤では、正面から当たる α 線だけを考えました。実際の α 線は、原子核から少し横にずれた所をねらって飛んできます。α 線がまっすぐ進んだとしたら原子核からいちばん近くを通る距離を、\textbf{ねらいの距離} $b$ といいます（専門書では衝突パラメータ）。\bkidx{ねらいのきょり}{ねらいの距離}$b$ が小さいほど原子核の近くを通り、大きく曲がります。ねらいの距離 $b$ と、曲がる角（\textbf{散乱角}）$\theta$ の関係を求めましょう。

\begin{bkfigure}
\begin{tikzpicture}[x=0.36mm, y=0.36mm]
  \draw[bkGray, line width=0.4pt, dash pattern=on 2pt off 1.5pt] (-150,0) -- (150,0);
  \draw[bkGray, line width=0.5pt, dash pattern=on 1pt off 1.5pt] (0,0) circle (45.5);
  \draw[txb path, -{Stealth[length=2mm]}] plot[smooth] coordinates {(-149.9,8.7) (-137.8,8.8) (-125.7,8.9) (-113.6,9.0) (-101.5,9.2) (-89.4,9.4) (-77.4,9.7) (-71.4,10.0) (-65.4,10.3) (-59.4,10.7) (-53.4,11.4) (-50.5,12.0) (-47.5,12.8) (-45.2,14.8) (-44.1,17.7) (-45.2,20.4) (-47.3,22.8) (-51.4,27.2) (-55.7,31.4) (-60.2,35.5) (-69.3,43.4) (-78.5,51.1) (-87.8,58.9) (-97.2,66.6) (-106.5,74.3) (-115.9,81.9)};
  \draw[txb path, -{Stealth[length=2mm]}] plot[smooth] coordinates {(-149.9,24.8) (-137.8,25.0) (-125.7,25.2) (-113.7,25.6) (-101.5,26.0) (-89.5,26.6) (-77.5,27.4) (-65.5,28.7) (-59.6,29.7) (-53.7,31.0) (-48.0,32.9) (-42.6,35.6) (-38.1,39.5) (-34.6,44.4) (-32.1,49.9) (-30.5,55.7) (-28.4,67.5) (-27.2,79.5) (-26.3,91.6) (-25.8,103.6) (-25.3,115.7) (-25.0,127.7)};
  \draw[txb path, -{Stealth[length=2mm]}] plot[smooth] coordinates {(-150.0,48.9) (-137.7,49.3) (-125.6,49.7) (-113.6,50.3) (-101.4,51.1) (-89.4,52.0) (-77.4,53.3) (-65.5,55.1) (-53.7,57.6) (-42.2,61.1) (-31.2,66.1) (-21.0,72.4) (-11.6,79.8) (-2.7,88.1) (5.7,96.7) (13.8,105.7) (21.6,114.9) (29.3,124.2)};
  \draw[txb path, -{Stealth[length=2mm]}] plot[smooth] coordinates {(-149.9,91.8) (-137.8,92.4) (-125.6,93.1) (-113.6,93.9) (-101.5,94.9) (-89.5,96.1) (-77.4,97.5) (-65.4,99.2) (-53.5,101.2) (-41.6,103.6) (-29.8,106.5) (-18.1,109.8) (-6.7,113.5) (4.7,117.6) (16.0,122.1) (27.2,126.8) (38.3,131.8) (49.2,137.0) (60.2,142.3) (71.0,147.7) (90,156.6) (110,165.9)};
  \draw[txb path, -{Stealth[length=2mm]}] (-150,1.6) -- (-48,1.6);
  \draw[txb path, -{Stealth[length=2mm]}] (-49,-2.6) -- (-120,-2.6);
  \txbnuc{(0,0)}{3.5}
  \node[txb lab, anchor=north] at (8,-6) {金の原子核};
  \node[txb lab, anchor=east] at (-153,10) {$b=\qty{8}{fm}$};
  \node[txb lab, anchor=east] at (-153,26) {\qty{23}{fm}};
  \node[txb lab, anchor=east] at (-153,49) {\qty{45}{fm}};
  \node[txb lab, anchor=east] at (-153,92) {\qty{85}{fm}};
  \node[txb lab, anchor=east] at (-153,-4) {正面};
  \node[txb lab, anchor=west] at (114,166) {$30^\circ$};
  \node[txb lab, anchor=west] at (34,124) {$60^\circ$};
  \node[txb lab, anchor=south] at (-25,131) {$90^\circ$};
  \node[txb lab, anchor=north east] at (-119,79) {$140^\circ$};
  \node[txb lab, anchor=north west, align=left] at (52,-6) {点線の円：正面から来たとき\\いちばん近づく距離 \qty{45}{fm}};
\end{tikzpicture}
\bkfigcaption{\qty{5.0}{MeV}の α 線が、金の原子核のそばを通るときの道すじ（この節の囲みのプログラムと同じ計算で描いた）。ねらいの距離 $b$ が小さいほど大きく曲がる。道すじは、原子核を焦点の1つとする双曲線になる。}\label{ph:原-1:fig-kidou}
\end{bkfigure}

\phboxhead{運動量の変わり方を2通りで数える}
α 線は、遠くから速さ $v$ で来て、遠くへ同じ速さ $v$ で去ります（クーロン力は保存力なので、遠くでの運動エネルギーは同じ。原子核の反動は無視）。向きだけが角 $\theta$ 変わります。運動量の変わり（変わったあと－前）の大きさは、2つの矢印 $m\vec v$ がつくる二等辺三角形の底辺で、
\[ |\Delta\vec p|=2mv\sin\frac{\theta}{2} \]
です。向きは、道すじの対称の軸（原子核と、いちばん近い点を結ぶ線）にそっています。

一方、運動量の変わりは、力を時間でたしたもの（力積：力-3）です。対称の軸の向きの力の成分は、α 線と原子核を結ぶ線が対称の軸となす角を $\varphi$ として $F\cos\varphi$、$F=k(2e)(Ze)/r^2$ です。ここで、原子核の周りの\textbf{角運動量}（力-6）$mr^2\,\dd\varphi/\dd t$ が、遠くでの $mvb$ と同じに保たれることを使うと、$\dd t=r^2\dd\varphi/(vb)$ となり、$r^2$ が消えて、
\[ |\Delta\vec p|=\int F\cos\varphi\,\dd t=\frac{2Zke^2}{vb}\int_{-\varphi_0}^{\varphi_0}\cos\varphi\,\dd\varphi=\frac{2Zke^2}{vb}\times2\sin\varphi_0 \]
です。$\varphi$ は、遠くから来るときの $-\varphi_0$ から、去るときの $+\varphi_0$ まで変わり、$\varphi_0=90^\circ-\theta/2$ なので、$\sin\varphi_0=\cos(\theta/2)$ です。2つの式を等しいとおき、$\tfrac12mv^2=K$、$d=2Zke^2/K$ を使うと、
\pheq{b=\frac{d}{2}\cot\frac{\theta}{2}=\frac{d}{2}\cdot\frac{1}{\tan(\theta/2)}}{m ＝ m}
\label{ph:原-1:cot}が出てきます（$\cot$ は $\tan$ の逆数：4-6⑤）。

\begin{phyomi}[$b=\dfrac{d}{2}\cot\dfrac{\theta}{2}$]
「$90^\circ$ 曲げるには、ねらいの距離を $d/2$ にする。それより大きく曲げるには、もっと真ん中をねらう」。$\theta=180^\circ$（真後ろにはね返る）は $b=0$（正面）、$\theta$ が小さいときは $b$ がとても大きい。ほんの少し曲がるだけの α 線は、原子核からずっと遠くを通ったものです。
\end{phyomi}

\begin{phbaisuru}[半分]{散乱角 $\theta$}
$\theta$ が小さい所では、$\cot(\theta/2)\approx2/\theta$（$\theta$ はラジアン）なので、ねらいの距離はおよそ2倍になります。小さく曲がる α 線ほど、遠くを通る「ねらいの広い」ものなので、数が多いのです。
\par\smallskip\noindent\begin{tikzpicture}
\begin{axis}[bkplot, width=100mm, height=42mm, xmin=0, xmax=182, ymin=0, ymax=300, xtick={10,30,60,90,120,150,180}, ytick={23,85,260}, xlabel={散乱角 $\theta$（度）}, ylabel={$b$（fm）}]
  \addplot[bkcurve, domain=8.8:180, samples=120] {22.75/tan(x/2)};
  \fill[bkRed] (axis cs:10,260) circle (1.8pt) (axis cs:30,84.9) circle (1.8pt) (axis cs:90,22.75) circle (1.8pt);
  \node[font=\sffamily\footnotesize, anchor=west, text=bkRed!75!black] at (axis cs:13,262) {$10^\circ$：\qty{260}{fm}};
  \node[font=\sffamily\footnotesize, anchor=south west, text=bkRed!75!black] at (axis cs:33,88) {$30^\circ$：\qty{85}{fm}};
  \node[font=\sffamily\footnotesize, anchor=south west, text=bkRed!75!black] at (axis cs:92,30) {$90^\circ$：\qty{23}{fm}（$d/2$）};
\end{axis}
\end{tikzpicture}
\end{phbaisuru}

\phboxhead{大きく曲がる α 線の割合}
はくの \qty{1}{m^2} に、原子核が $nt_{\mathrm f}$ 個並んでいるとします（$n$ は \qty{1}{m^3} あたりの原子の数、$t_{\mathrm f}$ ははくの厚さ）。1つの原子核の周りで、半径 $b$ の円の中をねらって来た α 線は、角 $\theta$ より大きく曲がります。はくがうすく、円どうしが重ならないとすると、$\theta$ より大きく曲がる α 線の割合は、円の面積の合計の割合で、
\pheq{(\theta\text{ より大きく曲がる割合})=nt_{\mathrm f}\times\pi b^2}{1 ＝ (1/m$^3$)×m×m$^2$}
です。\qty{0.40}{\micro m}の金はくで、\qty{5.0}{MeV}の α 線が $90^\circ$ より大きく曲がる割合は約 $3.8\times10^{-5}$、2万6千個に1個です（例題11）。ガイガーとマースデンが数えた「ごくまれ」は、この式で説明できました。円の面積 $\pi b^2$ のように、「1つの的の大きさ」で起こりやすさを表す量を、専門書では\textbf{断面積}とよびます（原-13 で詳しく使います）。

\Needspace{14\baselineskip}\phboxhead{プログラムで α 線を飛ばす}
\begin{phcode}[Python]
from math import sqrt, atan2, degrees
d = 45.5                  \# 正面から来たとき、いちばん近づく距離（fm）
for b in [200, 85, 40, 22.75, 10]:   \# ねらいの距離（fm）
    x, y = -20000.0, b    \# はじめの位置（fm）
    vx, vy = 1.0, 0.0     \# はじめの速さを1とする
    r = sqrt(x * x + y * y)
    while r < 20001:
        dt = r / 1000     \# 遠くでは大きく、近くでは小さく進める
        ax = d / 2 * x / r ** 3   \# 原子核から押される加速度
        ay = d / 2 * y / r ** 3
        vx = vx + ax * dt \# 速度を先に進める
        vy = vy + ay * dt
        x = x + vx * dt
        y = y + vy * dt
        r = sqrt(x * x + y * y)
    print(b, round(degrees(atan2(vy, vx)), 1))
\end{phcode}
\begin{phcodeout}[実行結果]
200 13.0
85 29.9
40 59.2
22.75 89.9
10 132.5
\end{phcodeout}
\noindent このプログラムは、α 線の動きを少しずつ進めるシミュレーションです（『力学』（力-16）の「速度を先に進める」方法）。速さを1、長さを fm で数えると、原子核から押される加速度の大きさは $(d/2)/r^2$ になります（$d$ の決め方 $\tfrac12mv^2=2Zke^2/d$ から）。原子核から遠い所では、力が弱いので大またに進め、近い所では細かく進めています。結果の2列めは曲がった角（度）で、$b=\qty{22.75}{fm}$（$d/2$）のとき、ほぼ $90^\circ$ です。式から出す値（13.0・30.0・59.3・90.0・132.5）と、0.1度ほどの違いで合っています。この計算で描いたのが、図\ref{ph:原-1:fig-kidou}の道すじです。

\bksubsection{9}{原子の模型の移りかわり}
\label{ph:原-1:sec9}
原子の中のようすは、実験が進むたびに書きかえられてきました。

\begin{bkfigure}
\begin{tikzpicture}[x=1mm, y=1mm]
  \foreach \cx/\ttl/\yr in {0/トムソン/1904年ごろ, 33/ラザフォード/1911年, 66/ボーア（原-4）/1913年, 99/電子の雲（原-7）/1926年から} {
    \node[txb lab, anchor=north, align=center] at (\cx,-14) {\ttl\\\yr};}
  \fill[bkRed!20] (0,0) circle (12); \draw[bkRed!60!black, line width=0.6pt] (0,0) circle (12);
  \foreach \p in {(-5,4),(4,6),(6,-3),(-3,-6),(0,0),(-8,-1)} {\txbel{\p}{1.3}}
  \draw[bkGray, line width=0.4pt] (33,0) circle (12);
  \fill[bkRed!75!black] (33,0) circle (0.8);
  \foreach \a in {30,150,260} {\txbel{($(33,0)+(\a:9)$)}{1.3}}
  \foreach \rr in {4,7.5,11} {\draw[bkGray, line width=0.5pt] (66,0) circle (\rr);}
  \fill[bkRed!75!black] (66,0) circle (0.8);
  \txbel{($(66,0)+(45:7.5)$)}{1.3}
  \shade[inner color=bkBlue!55, outer color=white] (99,0) circle (12);
  \fill[bkRed!75!black] (99,0) circle (0.8);
\end{tikzpicture}
\bkfigcaption{原子の模型の移りかわり。左から：正の電荷の球の中に電子がちらばる／小さな原子核の周りを電子が回る／電子は、とびとびの決まった軌道だけを回る／電子のいる場所は、確率の雲として広がる。}\label{ph:原-1:fig-mokei}
\end{bkfigure}

\noindent 1904年には、日本の長岡半太郎も、正の電荷をもつ重い芯の周りを、たくさんの電子が土星の輪のように回る模型を発表していました（東京での発表は1903年）。ラザフォードの模型は、この形に近いものです。ところが、ラザフォードの模型には大きな問題がありました。円を回る電子は加速度運動をしているので、電磁波（光）を出してエネルギーを失い（物理シリーズの『電磁気学』（電-15）で学んだ）、水素の原子なら、およそ $10^{-11}$ 秒（1000億分の1秒）ほどで原子核に落ちてしまうはずなのです。実際の原子はこわれません。この問題に答えたのがボーアの模型（原-4）で、さらに電子を波として扱う量子力学（第2部）で、原子のすがたは「電子の雲」として書かれるようになりました（原-7）。

\bksubsection{10}{電子の線を使う機械}
\label{ph:原-1:sec10}
陰極線（電子の線）を、電場と磁場で思いどおりに曲げる技術は、トムソンの実験から育ち、たくさんの機械になりました。
\begin{itemize}[leftmargin=1.2\zw, itemsep=2pt, topsep=2pt, label={\color{bkTheme}●}]
\item \textbf{ブラウン管}：①のテレビや、昔のオシロスコープ。電子を電圧で加速し、電場か磁場で曲げて画面に当てる。
\item \textbf{電子顕微鏡}：光の代わりに電子の線で物を見る。光の顕微鏡より何百倍も小さな物まで見える（理由は原-3）。
\item \textbf{X 線の管}：加速した電子を金属に当てて X 線を出す。病院のレントゲンや、空港の手荷物の検査（X 線の出方は原-4）。
\item \textbf{電子ビーム溶接・電子ビーム描画}：電子の線のエネルギーで金属を溶かしてつなぐ。半導体の細かい回路の型を描く。
\end{itemize}

\bksubsection{11}{質量分析：同位体を分ける}
\label{ph:原-1:sec11}
②の比電荷の考えを原子のイオンに使うと、原子の重さを見分けられます。原子から電子を1個はぎとった正のイオン（電荷 $+e$）を、電圧 $V$ で加速し、磁場 $B$ に垂直に入れます。加速で得る運動エネルギーは $eV$ なので、$\tfrac12mv^2=eV$ から $v=\sqrt{2eV/m}$、円の半径は
\pheq{r=\frac{mv}{eB}=\frac{1}{B}\sqrt{\frac{2mV}{e}}}{m ＝ (1/T)×$\sqrt{\text{kg·V/C}}$}
\label{ph:原-1:mass}です。半径は質量の平方根に比例するので、重いイオンほど大きな円をえがきます。半円を回ったところに検出器を並べると、質量ごとに別の場所に届きます。これが\textbf{質量分析計}です。\bkidx{しつりょうぶんせき}{質量分析}

\begin{bkfigure}
\begin{tikzpicture}[x=1mm, y=1mm]
  \fill[bkBlue!8] (-36,0) rectangle (36,36);
  \foreach \x in {-18,-6,6,18} {\foreach \y in {7,19} {\txbin{(\x,\y)}{1.5}}}
  \node[txb lab, anchor=south east] at (36,36.5) {磁場 $B$（紙のおくへ）};
  \draw[txb beam] (31.5,-14) -- (31.5,0);
  \draw[txb beam] (31.5,0) arc (0:180:31.5);
  \draw[bkPurple!80!black, line width=1.3pt, dash pattern=on 3pt off 1.5pt] (32.8,0) arc (0:180:32.8);
  \fill[txb plate] (29,-16) rectangle (34,-18); \node[txb lab, anchor=west] at (35,-17) {加速の電極（電圧 $V$）};
  \fill[bkGray!60] (-36,-4) rectangle (-28,-1.5);
  \node[txb lab, anchor=east] at (-37,-2.7) {検出器};
  \draw[txb dim] (-31.5,-9) -- (31.5,-9);
  \node[txb lab, anchor=north] at (0,-9.5) {直径 $2r$};
  \node[txb lab, anchor=west, align=left] at (40,20) {実線：$^{12}$C$^+$\\破線：$^{13}$C$^+$（少し大きい円）};
\end{tikzpicture}
\bkfigcaption{磁場で曲げる質量分析計（模式図）。イオンは下から入って半円を回る。重いイオンほど大きな半円をえがき、検出器の別の場所に届く。}\label{ph:原-1:fig-shitsuryou}
\end{bkfigure}

\begin{phmijika}[質量分析は、いろいろな所で]
質量分析は、同位体の割合を調べる道具として、炭素14 による年代の測定（原-12）や、岩石の年代の測定、病院の検査、食品やドーピングの検査、空気の中の物質の分析に使われます。今の質量分析計には、磁場で曲げる形のほかに、電場の中を飛ぶ時間をはかる形など、いろいろな形があります。どれも、比電荷で粒を見分けるという、トムソンの考えを受けついでいます。
\end{phmijika}

\bksubsection{12}{安全の約束：真空放電の管と X 線}
\label{ph:原-1:sec12}
\begin{column}{安全：真空放電の管は、先生が見せる}
この章の実験の多くは、高い電圧や、放射線を出す物を使います。どれも、自分でためすことはしません。
\begin{itemize}[leftmargin=1.2\zw, itemsep=1pt, topsep=2pt]
\item 真空放電の管（クルックス管など）は、誘導コイルで数万 V の高い電圧をかけます。電圧をかけた管からは、\textbf{X 線が出ることがある}ので、学校では先生が、電圧・距離・時間に気をつけて見せます。管に近づいたり、さわったりしないでください。
\item α 線を出す物質や霧箱は、科学館や大学で、専門の人がいる所で見ます。
\item 古いテレビ（ブラウン管）の中には、電源を切ったあとも高い電圧が残っていることがあります。分解しないでください。
\end{itemize}
この章で自分でためしてよいのは、電卓や表計算・プログラムを使った計算（やってみよう・⑧）です。
\end{column}

\begin{phmada}[陽子の大きさは？]
陽子の大きさ（電荷が広がっている半径）は、2つの方法ではかられてきました。1つは、電子を陽子に当てて、曲がり方を調べる方法（この章の α 線と同じ考え）。もう1つは、水素の原子が出す光の波長を、とても正確にはかる方法です（陽子が大きさをもつと、電子のエネルギーがわずかにずれる：原-4）。2010年、電子の代わりに、電子の約200倍重いミュー粒子を陽子の周りに回らせた「ミュー粒子の水素」の光をはかった実験で、それまでの値（約 \qty{0.88}{fm}）より4\,\% ほど小さい約 \qty{0.84}{fm} が出て、「陽子半径の謎」とよばれました。その後、2019年の電子を当てる新しい実験では \qty{0.831}{fm} が出るなど、小さいほうの値にまとまりつつあり、2022年にまとめられた値は \qty{0.84075}{fm} です。それでも、古い測定と食いちがった理由は完全にはわかっておらず、測定と検討が続いています。
\end{phmada}

\begin{phsaki}{電子の線と、原子をはかる機械}
この章では、一様な電場・磁場の中の動きだけを使いました。電子顕微鏡や加速器では、場所によって変わる電場・磁場で、電子の線をレンズのように集めたり広げたりします（電子光学）。光の速さに近い粒を加速する加速器では、相対論（原-11）を入れて設計します。質量分析は、分析化学や医学で、何万種類もの分子を見分ける道具に育っています。どれも、電子デバイス・加速器・分析化学の専門の本で学びます。
\end{phsaki}

\begin{senmonsho}[原子核物理の専門書の書き方]
\begin{bkhonbun}
ラザフォード散乱の微分断面積は
\[ \frac{\dd\sigma}{\dd\Omega}=\left(\frac{zZe^2}{16\pi\varepsilon_0K}\right)^2\frac{1}{\sin^4(\theta/2)} \]
で与えられる．
\end{bkhonbun}
\bkyaku{$z$ は飛んでくる粒の電荷の数（α 線は2）。かっこの中は $d/4$ と同じ大きさ（$1/(4\pi\varepsilon_0)=k$）。$\dd\Omega$ は立体角（8-8⑬）。「角 $\theta$ の向きの、決まった小さな見かけの大きさの所へ飛んでくる割合」を表す式で、⑧の「$\theta$ より大きく曲がる割合 $nt_{\mathrm f}\pi b^2$」を $\theta$ で微分して作る。$\sin^4(\theta/2)$ でわるので、角が小さいほど、ずっと多く飛んでくる。}
\end{senmonsho}

\begin{taisetsu}
\def\tsl#1{\makebox[7.6\zw][l]{\sffamily\bfseries\small\color{\bkTT}#1}}
\tsl{比電荷}$v=\dfrac{E}{B}$、$r=\dfrac{mv}{eB}$、$\dfrac{e}{m}=\dfrac{v}{rB}$。電子は $\qty{1.76e11}{C/kg}$\\[3pt]
\tsl{電気素量}$qE=mg$、$E=V/d_{\mathrm p}$。電荷は $e=\qty{1.60e-19}{C}$ の整数倍\\[3pt]
\tsl{終端速度}$\tfrac43\pi a^3\rho g=6\pi\eta av_{\mathrm t}$\\[3pt]
\tsl{近づく距離}$d=\dfrac{2Zke^2}{K}$、$ke^2=\qty{1.44}{MeV.fm}$\\[3pt]
\tsl{散乱}$b=\dfrac{d}{2}\cot\dfrac{\theta}{2}$、割合は $nt_{\mathrm f}\pi b^2$\\[3pt]
\tsl{原子核}陽子 $Z$ 個・中性子 $N$ 個、$A=Z+N$。半径 約 $1.2A^{1/3}$\,fm\\[3pt]
\tsl{原子の大きさ}約 \qtyrange{0.1}{0.3}{nm}。$N_{\mathrm A}=\num{96485}\,\mathrm{C/mol}\div e$\\[3pt]
\tsl{質量分析}$r=\dfrac{1}{B}\sqrt{\dfrac{2mV}{e}}$（半径は質量の平方根に比例）
\end{taisetsu}

\begin{kotoba}
\small
\bkword{陰極線}{真空放電の管の陰極から飛び出す電子の流れ。}
\bkword{比電荷}{粒の電荷÷質量。電場・磁場での曲がり方で決まる。}
\bkword{電気素量}{電荷の最小の単位 $e$。電子の電荷は $-e$、陽子は $+e$。}
\bkword{終端速度}{重さと抵抗がつり合って、一定になった落ちる速さ。}
\bkword{原子核}{原子のまん中の、正の電荷と質量のほとんどをもつ小さな部分。}
\bkword{陽子・中性子}{原子核をつくる粒。陽子は $+e$、中性子は電荷なし。まとめて核子。}
\bkword{原子番号・質量数}{陽子の数 $Z$、陽子と中性子の数の和 $A$。}
\bkword{同位体}{陽子の数が同じで、中性子の数がちがう原子。}
\bkword{ねらいの距離}{まっすぐ進んだとしたら、原子核からいちばん近くを通る距離 $b$。}
\bkword{質量分析}{イオンを磁場などで曲げ、質量で見分ける方法。}
\end{kotoba}

\bkkigocards{\bkkigocard{$e/m$}{イー・パー・エム}{比電荷〔C/kg〕。電荷÷質量。電子は $\qty{1.76e11}{C/kg}$}{charge-to-mass ratio}{原-1}}
{\bkkigocard{$Z$}{ゼット}{原子番号（陽子の数）。インピーダンス $Z$（波・電）、分配関数 $Z$（熱-15）とは別}{atomic number}{原-1}}
\bkkigocards{\bkkigocard{$b$}{ビー}{ねらいの距離（衝突パラメータ）〔m〕。波-14 のビット数 $b$ とは別}{impact parameter}{原-1}}
{\bkkigocard{$d$}{ディー}{正面から来た粒が、原子核にいちばん近づく距離〔m〕。原-1 では極板の間かくを $d_{\mathrm p}$ と書く}{distance of closest approach}{原-1}}

%% ======================== 例題 ========================
\Needspace{10\baselineskip}\bksection{reidai}{例題でたしかめよう}
\label{ph:原-1:ex1}
\begin{phreidai}[原子の中は、どのくらいからっぽ？]
金の原子核の半径を、目安の式 $R=1.2\times A^{1/3}$\,fm（$A=197$）で求めましょう。金の原子の半径 \qty{0.144}{nm} とくらべて、長さの比と体積の比を求めましょう。原子を直径 \qty{100}{m} の競技場にすると、原子核の直径は何 mm ですか。
\phbamen
はてな？の「原子の中はからっぽ」を、数で確かめる場面です。
\phmushi
原子も原子核も、はっきりした表面をもつ球とみます（本当は、はしがぼんやりしている）。
\phzu
図\ref{ph:原-1:fig-kyougijou}。原子核の半径 $R$、原子の半径 $r_{\text{原子}}=\qty{0.144}{nm}=\qty{1.44e-10}{m}$。
\phriyu
大きさの比は、同じ単位にそろえてわり算をする。体積は半径の3乗に比例するので、体積の比は長さの比の3乗（2-7）。
\phtateru
\[ R=1.2\times197^{1/3}\ \mathrm{fm},\qquad\text{長さの比}=\frac{R}{r_{\text{原子}}},\qquad\text{体積の比}=\left(\frac{R}{r_{\text{原子}}}\right)^3 \]
\phtoku
$197^{1/3}\approx5.82$ なので、$R\approx1.2\times5.82\approx\qty{7.0}{fm}=\qty{7.0e-15}{m}$。
\[ \frac{R}{r_{\text{原子}}}=\frac{7.0\times10^{-15}}{1.44\times10^{-10}}\approx4.8\times10^{-5},\qquad (4.8\times10^{-5})^3\approx1.1\times10^{-13} \]
競技場：$\qty{100}{m}\times4.8\times10^{-5}=\qty{4.8e-3}{m}$、約\qty{4.8}{mm}。
\phtashikame
単位：fm と m をそろえてからわったので、比は単位のない数。\\
大きさ：⑤で、α 線が近づいた距離 \qty{45}{fm} より小さく、つじつまが合う。体積で $10^{-13}$、つまり原子の体積の10兆分の1ほどしか原子核は占めていない。
\kotae 原子核の半径 約\qty{7.0}{fm}。長さの比 約 $4.8\times10^{-5}$、体積の比 約 $1.1\times10^{-13}$。競技場なら約\qty{4.8}{mm}
\end{phreidai}
\begin{jibun}
水素の原子では、陽子の半径は約 \qty{0.84}{fm}、原子の大きさの目安（ボーア半径：原-4）は \qty{0.0529}{nm} です。長さの比と体積の比を求めましょう。原子を直径 \qty{100}{m} にすると、陽子の直径は何 mm ですか。
\bkjans{長さの比 約 $1.6\times10^{-5}$、体積の比 約 $4.0\times10^{-15}$。陽子は約\qty{1.6}{mm}（金よりさらにからっぽ）}
\end{jibun}

\label{ph:原-1:ex2}
\begin{phreidai}[トムソンの方法で比電荷をはかる]
3けたで書きます。電子の線に、電場 $\qty{2.00e4}{V/m}$ と、それに垂直な磁場 \qty{1.00}{mT} をかけると、電子はまっすぐ進みました。電場を切ると、電子は半径 \qty{11.4}{cm} の円をえがきました。電子の速さと比電荷を求めましょう。
\phbamen
トムソンの実験で、電子の比電荷をはかる場面です（数は例のために置いた値とします）。
\phmushi
電子の重力、管に残った気体とのぶつかりを無視します。電子はどれも同じ速さとします。
\phzu
図\ref{ph:原-1:fig-thomson}。$E=\qty{2.00e4}{V/m}$、$B=\qty{1.00e-3}{T}$、$r=\qty{0.114}{m}$。
\phriyu
まっすぐ進むのは、電場の力と磁場の力がつり合うとき（$v=E/B$）。磁場だけでは等速円運動で、中心へ向かう力がローレンツ力（②）。
\phtateru
\[ eE=evB,\qquad m\frac{v^2}{r}=evB\qquad\text{より}\qquad v=\frac{E}{B},\quad \frac{e}{m}=\frac{v}{rB} \]
\phtoku
\[ v=\frac{2.00\times10^4}{1.00\times10^{-3}}=\qty{2.00e7}{m/s} \]
\[ \frac{e}{m}=\frac{2.00\times10^7}{0.114\times1.00\times10^{-3}}\approx\qty{1.75e11}{C/kg} \]
\phtashikame
単位：(m/s)÷(m×T)。T＝N/(A·m)＝kg/(C·s) なので、(m/s)÷(m·kg/(C·s))＝C/kg。\\
大きさ：今の値 $\qty{1.76e11}{C/kg}$ と 0.3\,\% で合う。速さは光の速さの約7\,\% で、相対論を考えなくてよい。
\kotae 速さ $\qty{2.00e7}{m/s}$、比電荷 約 $\qty{1.75e11}{C/kg}$
\end{phreidai}
\begin{jibun}
3けたで書きます。電場 $\qty{3.00e4}{V/m}$ と磁場 \qty{2.00}{mT} で電子がまっすぐ進み、電場を切ると半径 \qty{4.27}{cm} の円をえがきました。速さと比電荷を求めましょう。
\bkjans{速さ $\qty{1.50e7}{m/s}$。比電荷は $v/(rB)$ で、約 $\qty{1.76e11}{C/kg}$}
\end{jibun}

\label{ph:原-1:ex3}
\begin{phreidai}[電圧で加速した電子の円から比電荷を]
『電磁気学』（電-6 例題8）と同じ装置で、\qty{1.0}{kV}で加速した電子を、\qty{1.0}{mT}の磁場に垂直に入れました。円の半径をはかると \qty{10.66}{cm} でした。比電荷を求めましょう（はかった半径が4けたなので、答えも4けたで書きます）。
\phbamen
電-6 では、比電荷を知っていて半径を求めました。ここでは反対に、半径をはかって比電荷を求めます。
\phmushi
重力、相対論（速さは光の速さの約6\,\%：原-11）を無視します。はじめの速さは0とします。
\phzu
加速の電圧 $V=\qty{1000}{V}$、$B=\qty{1.0e-3}{T}$、$r=\qty{0.1066}{m}$。
\phriyu
加速で得るエネルギーは $eV$（電-2 ⑥）。円の半径は $r=mv/(eB)$（②）。2つの式から、はからない速さ $v$ を消す。
\phtateru
$\tfrac12mv^2=eV$ から $v^2=2eV/m$。$r=mv/(eB)$ を2乗すると $r^2=\dfrac{m^2v^2}{e^2B^2}=\dfrac{m}{e}\cdot\dfrac{2V}{B^2}$ なので、
\[ \frac{e}{m}=\frac{2V}{B^2r^2} \]
\phtoku
\[ \frac{e}{m}=\frac{2\times1000}{(1.0\times10^{-3})^2\times0.1066^2}=\frac{2000}{1.1364\times10^{-8}}\approx\qty{1.760e11}{C/kg} \]
\phtashikame
今の値 $\qty{1.7588e11}{C/kg}$ と 0.1\,\% で合う。\\
半径を電-6 のように2けたの \qty{0.11}{m} に丸めると、答えは $\qty{1.7e11}{C/kg}$ になり、6\,\% ずれる。2乗で入る量は、丸めのずれも2倍になる（$r$ が1\,\% ずれると、$e/m$ は約2\,\% ずれる）。
\kotae 約 $\qty{1.760e11}{C/kg}$
\end{phreidai}
\begin{jibun}
\qty{500}{V}で加速した電子を、\qty{1.0}{mT}の磁場に垂直に入れると、半径 \qty{7.54}{cm} の円をえがきました。比電荷を4けたで求めましょう。
\bkjans{$e/m=2V/(B^2r^2)$ で、約 $\qty{1.759e11}{C/kg}$}
\end{jibun}

\label{ph:原-1:ex4}
\begin{phreidai}[油滴を浮かべる電圧]
半径 \qty{1.0}{\micro m}、密度 \qty{870}{kg/m^3} の油滴があります。間かく \qty{5.0}{mm} の平らな板の間で、この油滴を浮かべるには、何 V の電圧が要りますか。油滴の電荷が $e$、$2e$、$3e$ のときで求めましょう。
\phbamen
ミリカンの実験で、浮かべる電圧の大きさを見当づける場面です（油の密度は例のための値とします）。
\phmushi
空気の浮力（油の密度の約700分の1）を無視します。板の間の電場は一様とします。
\phzu
図\ref{ph:原-1:fig-millikan}。$a=\qty{1.0e-6}{m}$、$\rho=\qty{870}{kg/m^3}$、$d_{\mathrm p}=\qty{5.0e-3}{m}$、$g=\qty{9.8}{m/s^2}$。
\phriyu
浮かんで止まっているので、電場の力と重さがつり合う（③）。電場は $E=V/d_{\mathrm p}$。
\phtateru
\[ q\,\frac{V}{d_{\mathrm p}}=mg\qquad\text{より}\qquad V=\frac{mgd_{\mathrm p}}{q},\qquad m=\frac43\pi a^3\rho \]
\phtoku
\[ m=\frac43\pi\times(1.0\times10^{-6})^3\times870\approx\qty{3.6e-15}{kg},\qquad mg\approx\qty{3.6e-14}{N} \]
\[ q=e:\ V=\frac{3.57\times10^{-14}\times5.0\times10^{-3}}{1.60\times10^{-19}}\approx\qty{1.1e3}{V} \]
$2e$ では半分の約\qty{560}{V}、$3e$ では3分の1の約\qty{370}{V}です。
\phtashikame
単位：N×m÷C＝J/C＝V。\\
大きさ：数百〜千 V で、実験の装置で出せる電圧。半径を \qty{10}{\micro m} にすると質量が1000倍になり、$e$ 1個では100万 V 以上要るので、小さな油滴を使う理由がわかる。
\kotae $e$：約 \qty{1.1e3}{V}、$2e$：約\qty{560}{V}、$3e$：約\qty{370}{V}
\end{phreidai}
\begin{jibun}
同じ油の、半径 \qty{0.80}{\micro m} の油滴を、間かく \qty{4.0}{mm} の板の間で浮かべます。油滴の電荷が $2e$ のとき、要る電圧を求めましょう。
\bkjans{$m\approx\qty{1.9e-15}{kg}$、$V=mgd_{\mathrm p}/(2e)\approx\qty{230}{V}$}
\end{jibun}

\label{ph:原-1:ex5}
\begin{phreidai}[ミリカンのデータから電気素量を]
いくつかの油滴の電荷をはかると、$3.2\times10^{-19}$、$4.8\times10^{-19}$、$8.0\times10^{-19}$、$11.2\times10^{-19}$\,C でした。電荷の最小の単位はいくらと考えられますか。
\phbamen
高校の入試でもよく出る、はかった電荷から $e$ を見つける場面です。
\phmushi
はかった値の誤差は、小さいとします（ここでは2けためまで正しい）。
\phzu
図\ref{ph:原-1:fig-seisu}のように、数直線に並べてみる。
\phriyu
どの電荷も、ある数 $e$ の整数倍のはず（③）。差をとると、$e$ の整数倍の小さな数が出てくる。
\phtateru
となりどうしの差：$4.8-3.2=1.6$、$8.0-4.8=3.2$、$11.2-8.0=3.2$（単位 $10^{-19}$\,C）。
\phtoku
差の中でいちばん小さいのは 1.6。どの電荷も 1.6 でわり切れるか確かめる。
\[ \frac{3.2}{1.6}=2,\quad\frac{4.8}{1.6}=3,\quad\frac{8.0}{1.6}=5,\quad\frac{11.2}{1.6}=7 \]
すべて整数になった。
\phtashikame
3.2 を単位とすると、$4.8/3.2=1.5$ で整数にならないので、3.2 は単位ではない。1.6 の半分の 0.8 も、すべてを整数倍にするが、0.8 の奇数倍の電荷が1つもないので、0.8 と考える理由はない（たくさんはかるほど、確かになる）。
\kotae 約 $\qty{1.6e-19}{C}$（電荷はその 2・3・5・7 倍）
\end{phreidai}
\begin{jibun}
はかった電荷が $6.4\times10^{-19}$、$9.6\times10^{-19}$、$14.4\times10^{-19}$、$17.6\times10^{-19}$\,C でした。電荷の最小の単位を求め、$3.2\times10^{-19}$\,C ではない理由も答えましょう。
\bkjans{$1.6\times10^{-19}$\,C（4・6・9・11 倍）。3.2 でわると $14.4/3.2=4.5$ などが整数にならないから}
\end{jibun}

\label{ph:原-1:ex6}
\begin{phreidai}[落ちる速さから油滴の大きさを]
電圧を切ると、油滴は一定の速さで落ち、顕微鏡の目もりで \qty{1.0}{mm} 落ちるのに \qty{9.6}{s} かかりました。油滴の半径を求めましょう。油の密度は \qty{870}{kg/m^3}、空気の粘度は $\qty{1.81e-5}{Pa.s}$ です。
\phbamen
ミリカンの方法で、はかりにのらない小さな油滴の大きさを、落ちる速さからはかる場面です。
\phmushi
空気の浮力を無視します。油滴は落ち始めてすぐに終端速度になるとします。
\phzu
$v_{\mathrm t}=\qty{1.0e-3}{m}\div\qty{9.6}{s}$、$\rho=870$、$\eta=\num{1.81e-5}$、$g=9.8$。
\phriyu
一定の速さで落ちるので、重さと空気の抵抗 $6\pi\eta av_{\mathrm t}$ がつり合っている（④）。
\phtateru
\[ \frac43\pi a^3\rho g=6\pi\eta av_{\mathrm t}\qquad\text{より}\qquad a=\sqrt{\frac{9\eta v_{\mathrm t}}{2\rho g}} \]
\phtoku
$v_{\mathrm t}\approx\qty{1.04e-4}{m/s}$ なので、
\[ a=\sqrt{\frac{9\times1.81\times10^{-5}\times1.04\times10^{-4}}{2\times870\times9.8}}=\sqrt{\frac{1.69\times10^{-8}}{\num{17052}}}\approx\qty{1.0e-6}{m} \]
\phtashikame
単位：(Pa·s×m/s)÷(kg/m$^3$×m/s$^2$)＝(N·s/m$^2$×m/s)÷(N/m$^3$)＝m$^2$ で、その平方根なので m。\\
大きさ：例題4の油滴と同じ大きさ。この油滴を浮かべる電圧をはかれば、例題4の式で電荷が出る。
\kotae 約 \qty{1.0}{\micro m}
\end{phreidai}
\begin{jibun}
同じ油の油滴が、\qty{2.0}{mm} 落ちるのに \qty{4.8}{s} かかりました。半径を求めましょう。
\bkjans{$v_{\mathrm t}\approx\qty{4.2e-4}{m/s}$、$a\approx\qty{2.0}{\micro m}$（速さが4倍で半径は2倍）}
\end{jibun}

\label{ph:原-1:ex7}
\begin{phreidai}[電気分解からアボガドロ定数を]
硝酸銀の水溶液に、\qty{0.500}{A}の電流を \qty{1930}{s} 流すと、−極に銀が \qty{1.079}{g} 出てきました。銀\qty{1}{mol}は\qty{107.87}{g}です。ファラデー定数とアボガドロ定数を求めましょう（3けたで書きます）。
\phbamen
電気分解の重さと、ミリカンの $e$ から、原子の数を数える場面です（電流と時間は例のための値とします）。
\phmushi
流れた電流は、すべて銀を出すのに使われたとします。
\phzu
電荷 $Q=It$、銀の物質量 $n_{\mathrm{Ag}}$、電子1個で銀1個（Ag$^+$ ＋ 電子 → Ag）。
\phriyu
銀 \qty{1}{mol} を出す電荷がファラデー定数。それは電子 $N_{\mathrm A}$ 個分の電荷なので、$e$ でわれば $N_{\mathrm A}$（⑦）。
\phtateru
\[ (\text{ファラデー定数})=\frac{Q}{n_{\mathrm{Ag}}},\qquad N_{\mathrm A}=\frac{(\text{ファラデー定数})}{e} \]
\phtoku
$Q=0.500\times1930=\qty{965}{C}$、$n_{\mathrm{Ag}}=1.079\div107.87\approx\qty{0.0100}{mol}$。
\[ \frac{965}{0.0100}\approx\qty{9.65e4}{C/mol},\qquad N_{\mathrm A}=\frac{9.65\times10^4}{1.60\times10^{-19}}\approx6.03\times10^{23}\ \mathrm{/mol} \]
\phtashikame
今の値 $6.02\times10^{23}$\,/mol と合う（ずれは $e$ を3けたに丸めたため）。\\
銅のイオン Cu$^{2+}$ なら、銅の原子1個に電子2個が要るので、\qty{1}{mol}を出す電荷は2倍になる。
\kotae ファラデー定数 約 $\qty{9.65e4}{C/mol}$、アボガドロ定数 約 $6.03\times10^{23}$\,/mol
\end{phreidai}
\begin{jibun}
硫酸銅の水溶液に \qty{1.00}{A} を \qty{3600}{s} 流すと、銅が \qty{1.185}{g} 出ました。銅\qty{1}{mol}は \qty{63.55}{g}、銅のイオンは Cu$^{2+}$ です。ファラデー定数とアボガドロ定数を3けたで求めましょう。
\bkjans{銅 \qty{0.01865}{mol}、電子はその2倍。ファラデー定数 約 $\qty{9.65e4}{C/mol}$、$N_{\mathrm A}$ 約 $6.03\times10^{23}$\,/mol}
\end{jibun}

\label{ph:原-1:ex8}
\begin{phreidai}[水の分子の大きさ]
水 \qty{1}{mol} は \qty{18.0}{g}、密度は \qty{1.00}{g/cm^3} です。水の分子1個が占める場所を立方体とみて、1辺の長さを求めましょう。
\phbamen
見えない分子の大きさを、ふだんはかれる量（重さ・体積）から見積もる場面です。
\phmushi
分子のすき間を無視し、分子が立方体のすき間なく並ぶとみます。
\phzu
\qty{1}{mol}の体積 $V_{\mathrm m}$、1個の体積 $V_{\mathrm m}/N_{\mathrm A}$、1辺 $l$。
\phriyu
\qty{1}{mol} の体積を、そこにある分子の数 $N_{\mathrm A}$ で等分する（⑦）。立方体の1辺は体積の3乗根。
\phtateru
\[ l=\left(\frac{V_{\mathrm m}}{N_{\mathrm A}}\right)^{1/3},\qquad V_{\mathrm m}=\frac{18.0\ \mathrm g}{1.00\ \mathrm{g/cm^3}} \]
\phtoku
$V_{\mathrm m}=\qty{18.0}{cm^3}=\qty{1.80e-5}{m^3}$。
\[ \frac{V_{\mathrm m}}{N_{\mathrm A}}=\frac{1.80\times10^{-5}}{6.02\times10^{23}}\approx\qty{2.99e-29}{m^3},\qquad l\approx\qty{3.1e-10}{m}=\qty{0.31}{nm} \]
\phtashikame
3乗根の見当：$3\times10^{-10}$ の3乗は $27\times10^{-30}=2.7\times10^{-29}$ で、近い。\\
大きさ：原子・分子の大きさ \qtyrange{0.1}{0.3}{nm} と合う。
\kotae 約 \qty{0.31}{nm}
\end{phreidai}
\begin{jibun}
金 \qty{1}{mol} は \qty{197}{g}、密度は \qty{19.3}{g/cm^3} です。金の原子1個が占める立方体の1辺を求めましょう。
\bkjans{1個の体積 約 $\qty{1.70e-29}{m^3}$、1辺 約\qty{0.26}{nm}（原子の半径 \qty{0.144}{nm} の直径 \qty{0.29}{nm} に近い）}
\end{jibun}

\label{ph:原-1:ex9}
\begin{phreidai}[α 線は原子核にどこまで近づくか]
\qty{5.0}{MeV}の α 線が、金の原子核（$Z=79$）に正面から向かいます。いちばん近づく距離を求めましょう。$ke^2=\qty{1.44}{MeV.fm}$ です。
\phbamen
ラザフォードの実験から、原子核の大きさの上限を見積もる場面です。
\phmushi
金の原子核の反動（α 線の約50倍重いので小さい）と、周りの電子の影響を無視します。原子核は点の電荷とみます。
\phzu
図\ref{ph:原-1:fig-kidou}の正面の道。α 線の電荷 $2e$、原子核 $Ze$、はじめの運動エネルギー $K$。
\phriyu
いちばん近づいた所で一瞬止まるので、運動エネルギーが全部クーロン力の位置エネルギーになる（エネルギーの保存：⑤）。
\phtateru
\[ K=\frac{k(2e)(Ze)}{d}\qquad\text{より}\qquad d=\frac{2Zke^2}{K} \]
\phtoku
\[ d=\frac{2\times79\times1.44\ \mathrm{MeV\cdot fm}}{5.0\ \mathrm{MeV}}\approx\qty{45}{fm} \]
\phtashikame
単位：MeV·fm÷MeV＝fm。\\
大きさ：原子の半径 \qty{0.144}{nm}（\num{144000}\,fm）の約3000分の1。α 線がここまで近づいてもクーロン力どおりに曲がったので、原子核はこれより小さい（今の値 約\qty{7}{fm}）。\\
極端な場合：$K$ を大きくすると $d$ は小さくなり、いつか原子核にふれる。
\kotae 約\qty{45}{fm}
\end{phreidai}
\begin{jibun}
\qty{7.7}{MeV}の α 線が、アルミニウムの原子核（$Z=13$、質量数27）に正面から向かうとき、いちばん近づく距離を求めましょう。アルミニウムの原子核の半径（$1.2\times27^{1/3}$\,fm）と α 線の半径（$1.2\times4^{1/3}$\,fm）の和とくらべると、何が言えますか。
\bkjans{$d=2\times13\times1.44\div7.7\approx\qty{4.9}{fm}$。半径の和は $3.6+1.9\approx\qty{5.5}{fm}$ で、$d$ のほうが小さい。α 線は原子核にふれるので、点の電荷のクーロン力だけのモデルはくずれる（実際、アルミニウムでは式からずれた散乱が見つかった）}
\end{jibun}

\label{ph:原-1:ex10}
\begin{phreidai}[曲がる角と、ねらいの距離]
例題9の α 線（$d=\qty{45.5}{fm}$）が、$5^\circ$・$30^\circ$・$90^\circ$・$150^\circ$ 曲がるときの、ねらいの距離を求めましょう。
\phbamen
大きく曲がる α 線が、原子核のどれくらい近くをねらって来たものかを調べる場面です。
\phmushi
原子核の反動と、周りの電子を無視します。1つの原子核だけで曲がるとします。
\phzu
図\ref{ph:原-1:fig-kidou}。ねらいの距離 $b$、散乱角 $\theta$。
\phriyu
⑧の式 $b=(d/2)\cot(\theta/2)$。電卓では $\cot$ を $1\div\tan$ で計算する（DEG にする：4-6）。
\phtateru
\[ b=\frac{45.5}{2}\times\frac{1}{\tan(\theta/2)}=\frac{22.75}{\tan(\theta/2)}\ \mathrm{fm} \]
\phtoku
\[ 5^\circ:\ \frac{22.75}{\tan2.5^\circ}=\frac{22.75}{0.0437}\approx\qty{521}{fm},\qquad 30^\circ:\ \frac{22.75}{\tan15^\circ}=\frac{22.75}{0.268}\approx\qty{85}{fm} \]
\[ 90^\circ:\ \frac{22.75}{\tan45^\circ}=\qty{22.75}{fm},\qquad 150^\circ:\ \frac{22.75}{\tan75^\circ}=\frac{22.75}{3.73}\approx\qty{6.1}{fm} \]
\phtashikame
⑧のプログラムで $b=85$ と 22.75 を入れると、$29.9^\circ$ と $89.9^\circ$ になり合う。\\
大きさ：$150^\circ$ に曲がるには、原子核の半径（\qty{7}{fm}）くらいの所をねらう必要がある。はく全体の面積のうち、そこはほんのわずか。
\kotae $5^\circ$：約\qty{521}{fm}、$30^\circ$：約\qty{85}{fm}、$90^\circ$：約\qty{23}{fm}、$150^\circ$：約\qty{6.1}{fm}
\end{phreidai}
\begin{jibun}
\qty{7.7}{MeV}の α 線と金の原子核では、$d\approx\qty{29.5}{fm}$ です。$60^\circ$ と $90^\circ$ 曲がるときのねらいの距離を求めましょう。
\bkjans{$60^\circ$：$14.77\div\tan30^\circ\approx\qty{26}{fm}$、$90^\circ$：約\qty{15}{fm}}
\end{jibun}

\label{ph:原-1:ex11}
\begin{phreidai}[大きく曲がる α 線の割合]
厚さ \qty{0.40}{\micro m} の金はくに、\qty{5.0}{MeV}の α 線を当てます。$90^\circ$ より大きく曲がる α 線の割合を求めましょう。金の密度は \qty{19.32}{g/cm^3}、\qty{1}{mol}は\qty{196.97}{g}です。
\phbamen
ガイガーとマースデンが数えた「ごくまれ」を、ラザフォードの模型で計算する場面です（はくの厚さは例のための値とします）。
\phmushi
1つの α 線は、いちばん近くを通った1つの原子核だけで曲がるとします（何回も少しずつ曲がることを無視）。原子核の反動と電子を無視します。
\phzu
数密度 $n$、厚さ $t_{\mathrm f}=\qty{4.0e-7}{m}$、$90^\circ$ のねらいの距離 $b_{90}=d/2=\qty{22.75}{fm}$。
\phriyu
$b<b_{90}$ をねらった α 線だけが $90^\circ$ より大きく曲がる。その円の面積の合計の割合が答え（⑧）。
\phtateru
\[ n=\frac{19.32\times10^3\ \mathrm{kg/m^3}}{196.97\times10^{-3}\ \mathrm{kg/mol}}\times N_{\mathrm A},\qquad(\text{割合})=nt_{\mathrm f}\times\pi b_{90}^2 \]
\phtoku
\qty{1}{m^3} あたり $n\approx5.90\times10^{28}$ 個、\qty{1}{m^2} あたり $nt_{\mathrm f}\approx2.36\times10^{22}$ 個。
\[ \pi b_{90}^2=\pi\times(2.275\times10^{-14})^2\approx\qty{1.63e-27}{m^2} \]
\[ (\text{割合})\approx2.36\times10^{22}\times1.63\times10^{-27}\approx3.8\times10^{-5} \]
\phtashikame
逆数は約 \num{26000}。2万6千個に1個が $90^\circ$ より大きく曲がる。\\
円が重ならないか：\qty{1}{m^2} に $2.36\times10^{22}$ 個なので、原子核1個あたりの面積は $\qty{4.2e-23}{m^2}$。円の面積はその約2万6千分の1で、重ならない。
\kotae 約 $3.8\times10^{-5}$（約2万6千個に1個）
\end{phreidai}
\begin{jibun}
例題11と同じはくと α 線で、$30^\circ$ より大きく曲がる割合を求めましょう（例題10の $b$ を使う）。
\bkjans{$b_{30}=\qty{85}{fm}$ で、割合は約 $5.3\times10^{-4}$（約1900個に1個。$90^\circ$ より大きいものの約14倍）}
\end{jibun}

\label{ph:原-1:ex12}
\begin{phreidai}[質量分析で炭素の同位体を分ける]
炭素12 と炭素13 の1価の正のイオン（$^{12}$C$^+$、$^{13}$C$^+$）を、\qty{1.0}{kV}で加速し、\qty{0.10}{T}の磁場に垂直に入れます。それぞれの円の半径と、半円を回ったあとに届く場所の差を求めましょう。$^{12}$C$^+$ の質量は $12\times\qty{1.66e-27}{kg}$、$^{13}$C$^+$ は $13.0\times\qty{1.66e-27}{kg}$ とします。
\phbamen
質量分析計で同位体を見分ける場面です。
\phmushi
イオンのはじめの速さ、電子1個分の質量、重力を無視します。
\phzu
図\ref{ph:原-1:fig-shitsuryou}。$V=\qty{1000}{V}$、$B=\qty{0.10}{T}$、電荷 $e$。
\phriyu
加速のエネルギー $eV$ と、円の半径 $r=mv/(eB)$（⑪）。届く場所は、入口から直径 $2r$ の所。
\phtateru
\[ r=\frac{1}{B}\sqrt{\frac{2mV}{e}} \]
\phtoku
\[ ^{12}\mathrm C^+:\ r=\frac{1}{0.10}\sqrt{\frac{2\times12\times1.66\times10^{-27}\times1000}{1.60\times10^{-19}}}\approx\qty{0.158}{m} \]
$^{13}$C$^+$ の半径は、質量の平方根に比例するので
\[ ^{13}\mathrm C^+:\ r\approx0.158\times\sqrt{\frac{13.0}{12}}\approx0.158\times1.041\approx\qty{0.164}{m} \]
届く場所の差は、直径の差なので $2\times(0.164-0.158)\approx\qty{0.013}{m}$。
\phtashikame
単位：V＝J/C、J＝kg·m$^2$/s$^2$ を使うと、根号の中は kg$^2$·m$^2$/(C$^2$·s$^2$) で、根号をとると kg·m/(C·s)。これを T＝kg/(C·s) でわって m。\\
大きさ：質量が $1/12$（約8\,\%）ちがうと、半径は約4\,\% ちがう（平方根なので半分）。
\kotae $^{12}$C$^+$：約\qty{0.158}{m}、$^{13}$C$^+$：約\qty{0.164}{m}。届く場所の差 約\qty{1.3}{cm}
\end{phreidai}
\begin{jibun}
塩素の同位体のイオン $^{35}$Cl$^+$ と $^{37}$Cl$^+$ を、\qty{2.0}{kV}で加速し、\qty{0.20}{T}の磁場に垂直に入れます。質量は、それぞれ $35.0\times1.66\times10^{-27}$\,kg と $37.0\times1.66\times10^{-27}$\,kg とします。半径と、届く場所の差を求めましょう。
\bkjans{$^{35}$Cl$^+$：約\qty{0.191}{m}、$^{37}$Cl$^+$：約\qty{0.196}{m}、差 約\qty{1.1}{cm}}
\end{jibun}

\begin{column}{電子は粒、そして波：トムソンの父と子}
1897年に陰極線が「粒」の流れだと示したのは、J. J. トムソンでした。それから30年ほどあと、息子の G. P. トムソンは、電子の線をうすい金属のはくに通すと、波のように回折することを示しました。父は電子が粒であることで（1906年）、子は電子が波であることで（1937年、デビソンと共に）、それぞれノーベル物理学賞を受けています。電子は粒なのか、波なのか。その答えは原-3 で考えます。
\end{column}

%% ======================== まちがえやすいところ ========================
\Needspace{13\baselineskip}\bksection{tsumazuki}{まちがえやすいところをチェック}
\begin{tsumazuki}{原子核は、原子の半分くらいの大きさ？}
\bkng{原子の絵では、原子核が原子の半分くらいの大きさに描いてある。}
\bkok{原子核の直径は、原子の約2万分の1。絵は見やすくするために、原子核をとても大きく描いている。}
本当の大きさで描くと、原子核は点にしか見えません（図\ref{ph:原-1:fig-kyougijou}）。
\end{tsumazuki}

\begin{tsumazuki}{比電荷がわかれば、電荷も質量もわかる？}
\bkng{トムソンの実験で、電子の電荷と質量が別々にわかった。}
\bkok{わかったのは比 $e/m$ だけ。$e$ はミリカンの実験ではかり、$m=e\div(e/m)$ で質量が出た。}
電場・磁場の中の動きは、加速度 $(q/m)(\vec E+\vec v\times\vec B)$ で決まるので、$q$ と $m$ は比でしか出てきません（②）。
\end{tsumazuki}

\begin{tsumazuki}{電子が受ける磁場の力の向き}
\bkng{電子が受けるローレンツ力の向きは、フレミングの左手の法則の親指の向き。}
\bkok{左手の法則の中指は「電流の向き」。電子は負なので、電子の動く向きと電流の向きは反対で、力も反対向きになる。}
図\ref{ph:原-1:fig-thomson}で、正の電荷なら上向きの力が、電子では下向きになります（②）。
\end{tsumazuki}

\begin{tsumazuki}{ねらいの距離は、いちばん近づく距離？}
\bkng{ねらいの距離 $b$ の所まで、α 線は原子核に近づく。}
\bkok{$b$ は「まっすぐ進んだとしたら」の距離。退け合う力で道が曲がるので、本当にいちばん近づく距離は $b$ より大きい。}
図\ref{ph:原-1:fig-kidou}の $b=\qty{23}{fm}$ の道は、原子核から \qty{23}{fm} より遠い所で曲がっています（⑧）。
\end{tsumazuki}

\begin{tsumazuki}{質量分析の半径は、質量に比例する？}
\bkng{質量が2倍のイオンは、半径も2倍。}
\bkok{同じ電圧で加速したイオンでは、半径は質量の\textbf{平方根}に比例する。2倍なら約1.41倍。}
重いイオンは、同じ電圧で加速しても遅いので、$r=mv/(eB)$ の $mv$ は $\sqrt m$ に比例します（⑪）。同じ速さのイオンなら、半径は質量に比例します。
\end{tsumazuki}


%% ======================== 工学ではこう使う ========================
\Needspace{10\baselineskip}\bksection{phgear}{工学ではこう使う}
\begin{phkougaku}[電子の線と、イオンの線]
\textbf{電子の線を曲げる機械}：電子顕微鏡、X 線の管、電子ビーム溶接、半導体の回路の型を描く電子ビーム描画装置は、どれも、②の比電荷の考えで電子の道を設計します。電子は比電荷がとても大きいので、\qty{1}{mT}ほどの弱い磁場でも、半径 \qty{10}{cm} ほどの円に曲がります（例題3）。

\textbf{イオンの線を曲げる機械}：質量分析計（⑪）や、半導体にホウ素やリンを打ちこむイオン注入装置（原-9）は、イオンの質量の違いで道が変わることを使います。イオンの比電荷は電子の数千〜数万分の1なので、同じ半径に曲げるには、ずっと強い磁場（\qty{0.1}{T}以上）が要ります。

\textbf{小さな粒の電荷}：コピー機のトナー、電気集じん器（電-1）、静電塗装では、粒の電荷と重さ・空気の抵抗のつり合いを、ミリカンの油滴と同じ考えで設計します。

\begin{phketa}[原子と電子の数の目安]
\phketaline{電子の比電荷・陽子の比電荷}{$\qty{1.76e11}{C/kg}$・$\qty{9.58e7}{C/kg}$}
\phketaline{電気素量}{$e=\qty{1.60e-19}{C}$（\qty{1}{C}は電子 $6.25\times10^{18}$ 個分）}
\phketaline{原子・小さな分子の大きさ}{\qtyrange{0.1}{0.3}{nm}}
\phketaline{原子核の半径}{\qtyrange{1}{8}{fm}（$1.2A^{1/3}$\,fm）}
\phketaline{クーロンの位置エネルギーの目安}{$ke^2=\qty{1.44}{MeV.fm}=\qty{1.44}{eV.nm}$}
\phboxhead{答えが桁で合っているか}
原子の大きさの答えが \qty{1}{nm} より大きく出たり、原子核の大きさが \qty{100}{fm} より大きく出たりしたら、単位の換算（nm＝$10^{-9}$\,m、fm＝$10^{-15}$\,m、MeV＝$\qty{1.60e-13}{J}$）を見直しましょう。
\end{phketa}
\end{phkougaku}
\phstation{電子デバイス,真空工学,分析化学,加速器工学}{比電荷、電子銃、電子光学、電気素量、質量分析計、同位体比、衝突パラメータ、散乱断面積}

\begin{column}{今は、原子1個を「見る」ことができる}
\textbf{走査トンネル顕微鏡}は、とがった針を表面から \qty{1}{nm} ほどまで近づけ、間に流れるわずかな電流で、表面の凹凸を原子1個ずつなぞって画像にします（発明者は1986年にノーベル物理学賞）。すき間を電流が流れるしくみは原-6 で学びます。
\end{column}

%% ======================== 章末問題 ========================
\bkendproblems
\bkproblemgroup{基本}{3}
\bkproblem 電場 $\qty{1.5e4}{V/m}$ と、それに垂直な磁場 \qty{1.0}{mT} をかけると、電子はまっすぐ進みました。電場を切ると、半径 \qty{8.5}{cm} の円をえがきました。（1）電子の速さを求めましょう。（2）比電荷を求めましょう。
\bkpans{（1）$v=E/B=\qty{1.5e7}{m/s}$。（2）$e/m=v/(rB)$ で、約 $\qty{1.8e11}{C/kg}$}
\bkproblem 油滴の電荷をはかると、$4.8\times10^{-19}$、$8.0\times10^{-19}$、$11.2\times10^{-19}$、$14.4\times10^{-19}$\,C でした。電気素量を求めましょう。
\bkpans{差は 3.2（単位 $10^{-19}$\,C）。$3.2$ では $4.8/3.2=1.5$ となり合わない。$1.6\times10^{-19}$\,C でわると 3・5・7・9 で、すべて整数。$\qty{1.6e-19}{C}$}
\bkproblem 半径 \qty{0.90}{\micro m}、密度 \qty{870}{kg/m^3} の油滴を、間かく \qty{6.0}{mm} の板の間で浮かべます。電荷が $3e$ のとき、要る電圧を求めましょう。
\bkpans{$m=\tfrac43\pi a^3\rho\approx\qty{2.7e-15}{kg}$。$V=mgd_{\mathrm p}/(3e)\approx\qty{330}{V}$}
\bkproblem \qty{5.0}{MeV}の α 線が、銀の原子核（$Z=47$）に正面から向かいます。（1）いちばん近づく距離を求めましょう。（2）α 線のエネルギーを \qty{10}{MeV} にすると、いくらになりますか。
\bkpans{（1）$2\times47\times1.44\div5.0\approx\qty{27}{fm}$。（2）$K$ に反比例するので半分の約\qty{14}{fm}}
\bkproblem 水素のイオン $^1$H$^+$ と、重水素のイオン $^2$H$^+$ を、\qty{1.0}{kV}で加速し、\qty{0.10}{T}の磁場に垂直に入れます。質量は、それぞれ $1.67\times10^{-27}$\,kg と $3.34\times10^{-27}$\,kg とします。それぞれの半径を求めましょう。半径の比はいくらですか。
\bkpans{$r=\sqrt{2mV/e}\div B$ で、$^1$H$^+$ は約\qty{4.6}{cm}、$^2$H$^+$ は約\qty{6.5}{cm}。比は $\sqrt2\approx1.4$}
\bkproblem 鉄 \qty{1}{mol} は \qty{55.85}{g}、密度は \qty{7.874}{g/cm^3} です。鉄の原子1個が占める立方体の1辺を求めましょう。
\bkpans{1個の体積 約 $\qty{1.18e-29}{m^3}$、1辺 約\qty{0.23}{nm}}
\bkproblem アルミニウム（$A=27$）の原子核の半径を $1.2\times A^{1/3}$\,fm で求め、原子の半径 \qty{0.143}{nm} との長さの比を求めましょう。原子を直径 \qty{1.0}{km} にすると、原子核の直径は何 cm ですか。
\bkpans{$27^{1/3}=3$ なので約\qty{3.6}{fm}。長さの比 約 $2.5\times10^{-5}$。原子核の直径は約\qty{2.5}{cm}}
\bkproblemgroup{チャレンジ}{4}
\bkproblem 例題11の金はくを2倍の厚さにし、α 線のエネルギーも2倍の \qty{10}{MeV} にしました。$90^\circ$ より大きく曲がる α 線の割合は、例題11の何倍になりますか。何個に1個になりますか。
\bkpans{割合は $nt_{\mathrm f}\pi b_{90}^2$、$b_{90}=d/2\propto1/K$ なので、厚さで2倍、$b_{90}^2$ で $\tfrac14$ 倍。あわせて $\tfrac12$ 倍の約 $1.9\times10^{-5}$、約5万2千個に1個}
\bkproblem 電圧を切ると、ある油滴は \qty{1.0}{mm} 落ちるのに \qty{15.0}{s} かかりました。間かく \qty{5.0}{mm} の板に \qty{190}{V} をかけると、この油滴は浮かんで止まりました。油の密度は \qty{870}{kg/m^3}、空気の粘度は $\qty{1.81e-5}{Pa.s}$ です。油滴の半径・質量・電荷を求め、電荷は $e$ の何倍かを答えましょう。
\bkpans{$v_{\mathrm t}\approx\qty{6.7e-5}{m/s}$、$a=\sqrt{9\eta v_{\mathrm t}/(2\rho g)}\approx\qty{0.80}{\micro m}$、$m\approx\qty{1.85e-15}{kg}$、$q=mgd_{\mathrm p}/V\approx\qty{4.8e-19}{C}$。$e$ の3倍}
\bkproblem 例題11の金はくで、散乱角が $30^\circ$ から $90^\circ$ までの α 線の割合を求めましょう。また、$\theta$ より大きく曲がる割合が $\cot^2(\theta/2)$ に比例することを示し、$30^\circ$ と $90^\circ$ の場合の比を求めましょう。
\bkpans{割合は $nt_{\mathrm f}\pi(b_{30}^2-b_{90}^2)$。$b_{30}=\qty{84.9}{fm}$、$b_{90}=\qty{22.75}{fm}$ で、約 $5.0\times10^{-4}$。割合 $nt_{\mathrm f}\pi b^2$ に $b=(d/2)\cot(\theta/2)$ を入れると $\cot^2(\theta/2)$ に比例する。比は $\cot^215^\circ\div\cot^245^\circ\approx13.9$}

\begin{furikaeri}
\begin{checklist}
\item 陰極線が負の電荷の粒（電子）の流れであることを、実験の結果から説明できる。\bkfurref{ph:原-1:sec1}
\item トムソンの方法で、電子の比電荷を求められる。\bkfurref{ph:原-1:sec2}
\item ミリカンの方法で、油滴の大きさと電荷を求め、電気素量を見つけられる。\bkfurref{ph:原-1:sec3}
\item ラザフォードの実験から、原子核が小さく重いことを、いちばん近づく距離で説明できる。\bkfurref{ph:原-1:sec5}
\item ねらいの距離と散乱角の式で、大きく曲がる α 線の割合を計算できる。\bkfurref{ph:原-1:sec8}
\item 質量分析で、同位体のイオンが届く場所の違いを求められる。\bkfurref{ph:原-1:sec11}
\end{checklist}
\end{furikaeri}
\bkfukidashi{sora}{机も体も、ほとんどからっぽの原子でできているなんて！ それなのに手が机を通りぬけないのは、どうしてだろう？}
